\documentclass[UTF8, 12pt, fontset=fandol, oneside]{ctexart}
\usepackage{researchnotes}

\title{第九章\quad 内积空间}
\author{}
\date{}

\begin{document}

\maketitle

\section*{§9.1\quad 内积空间的概念}

在解析几何中，我们已经知道，$\mathbb R^3$ 中任一向量可定义其``长度''，空间中任意两点之间可定义其距离. 向量 $v=(x_1,x_2,x_3)$ 的长度为
\[
  \|v\|=\sqrt{x_1^2+x_2^2+x_3^2}.
\]
两点 $(x_1,x_2,x_3),(y_1,y_2,y_3)$ 之间的距离为
\[
  \sqrt{(x_1-y_1)^2+(x_2-y_2)^2+(x_3-y_3)^2},
\]
即这两点所代表的向量之差的长度. 我们现在要把``长度''、``距离''的概念推广到一般的实线性空间与复线性空间上去. 距离可看成是长度的派生概念，而长度又可看成是内积的派生概念. 我们已经知道在 $\mathbb R^3$ 中，内积是这样定义的：若 $u=(x_1,x_2,x_3),v=(y_1,y_2,y_3)$，则 $u$ 与 $v$ 的内积（或点积）为
\begin{equation}
  u\cdot v=x_1y_1+x_2y_2+x_3y_3,
  \tag{9.1.1}
\end{equation}
从而 $u$ 的长度为
\[
  \|u\|=(u\cdot u)^{\frac12}.
\]
从 (9.1.1) 式我们可看出 $\mathbb R^3$ 中的内积有下列性质.

$\mathbb R^3$ 中向量内积的性质
\[
\begin{gathered}
  (1)\quad u\cdot v=v\cdot u;\\
  (2)\quad (u+w)\cdot v=u\cdot v+w\cdot v;\\
  (3)\quad (cu)\cdot v=c(u\cdot v);\\
  (4)\quad u\cdot u\ge 0,\text{ 且 }u\cdot u=0\text{ 当且仅当 }u=0,
\end{gathered}
\]
其中 $u,v,w$ 是 $\mathbb R^3$ 中的任意向量，$c$ 是任一实数.

根据上述 4 条性质，我们类似地定义一般线性空间上的内积：

\noindent\textbf{定义 9.1.1}\quad 设 $V$ 是实数域上的线性空间，若存在某种规则，使对 $V$ 中任意一组有序向量 $\{\alpha,\beta\}$，都唯一地对应一个实数，记为 $(\alpha,\beta)$，且适合如下规则：
\[
\begin{gathered}
  (1)\quad (\beta,\alpha)=(\alpha,\beta);\\
  (2)\quad (\alpha+\beta,\gamma)=(\alpha,\gamma)+(\beta,\gamma);\\
  (3)\quad (c\alpha,\beta)=c(\alpha,\beta),\ c\text{ 为任一实数};\\
  (4)\quad (\alpha,\alpha)\ge0\text{ 且等号成立当且仅当 }\alpha=0,
\end{gathered}
\]
则称在 $V$ 上定义了一个内积. 实数 $(\alpha,\beta)$ 称为 $\alpha$ 与 $\beta$ 的内积. 线性空间 $V$ 称为实内积空间. 有限维实内积空间称为 Euclid 空间，简称为欧氏空间.

对复数域上的线性空间，我们也可以定义内积.

\noindent\textbf{定义 9.1.2}\quad 设 $V$ 是复数域上的线性空间，若存在某种规则，使对 $V$ 中任意一组有序向量 $\{\alpha,\beta\}$，都唯一地对应一个复数，记为 $(\alpha,\beta)$，且适合如下规则：
\[
\begin{gathered}
  (1)\quad (\beta,\alpha)=\overline{(\alpha,\beta)};\\
  (2)\quad (\alpha+\beta,\gamma)=(\alpha,\gamma)+(\beta,\gamma);\\
  (3)\quad (c\alpha,\beta)=c(\alpha,\beta),\ c\text{ 为任一复数};\\
  (4)\quad (\alpha,\alpha)\ge0\text{ 且等号成立当且仅当 }\alpha=0,
\end{gathered}
\]
则称在 $V$ 上定义了一个内积. 复数 $(\alpha,\beta)$ 称为 $\alpha$ 与 $\beta$ 的内积. 线性空间 $V$ 称为复内积空间. 有限维复内积空间称为酉空间.

\noindent\textbf{注}\quad 实内积空间的定义与复内积空间的定义是相容的. 事实上，对一个实数 $a$，$\overline a=a$，故定义 9.1.1 中的 (1) 与定义 9.1.2 中的 (1) 是一致的. 因此，我们经常将这两种空间统称为内积空间，在某些定理的叙述及证明中也不区分它们，而统一作为复内积空间来处理. 但是，需要注意的是对复内积空间，定义 9.1.2 中的 (1),(3) 意味着：
\[
  (\alpha,c\beta)=\overline c(\alpha,\beta).
\]

\noindent\textbf{例 9.1.1}\quad 设 $\mathbb R_n$ 是 $n$ 维实列向量空间，$\alpha=(x_1,x_2,\cdots,x_n)'$，$\beta=(y_1,y_2,\cdots,y_n)'$，定义
\[
  (\alpha,\beta)=x_1y_1+x_2y_2+\cdots+x_ny_n,
\]
则在此定义下 $\mathbb R_n$ 成为一个欧氏空间，上述内积称为 $\mathbb R_n$ 的标准内积.

\noindent\textbf{例 9.1.2}\quad 设 $\mathbb C_n$ 是 $n$ 维复列向量空间，$\alpha=(x_1,x_2,\cdots,x_n)'$，$\beta=(y_1,y_2,\cdots,y_n)'$，定义
\[
  (\alpha,\beta)=x_1\overline{y_1}+x_2\overline{y_2}+\cdots+x_n\overline{y_n},
\]
则在此定义下 $\mathbb C_n$ 成为一个酉空间，上述内积称为 $\mathbb C_n$ 的标准内积.

\noindent\textbf{注}\quad 对 $n$ 维实或复行向量空间，我们也可同样定义标准内积.

\noindent\textbf{例 9.1.3}\quad 设 $V=\mathbb R^2$ 为二维实行向量空间，若 $\alpha=(x_1,x_2),\beta=(y_1,y_2)$，定义
\[
  (\alpha,\beta)=x_1y_1-x_2y_1-x_1y_2+4x_2y_2,
\]
容易验证定义 9.1.1 中的 (1),(2),(3) 都成立. 当 $\beta=\alpha$ 时，上式为
\[
  (\alpha,\alpha)=x_1^2-2x_1x_2+4x_2^2=(x_1-x_2)^2+3x_2^2,
\]
故 (4) 也成立. 因此，$\mathbb R^2$ 在此内积下成为二维欧氏空间.

\noindent\textbf{例 9.1.4}\quad 设 $V$ 是由 $[a,b]$ 区间上连续函数全体构成的实线性空间，设 $f(t),g(t)\in V$，定义
\[
  (f,g)=\int_a^b f(t)g(t)\mathrm dt,
\]
则不难验证这是一个内积，于是 $V$ 成为内积空间. 这是一个无限维实内积空间.

\noindent\textbf{例 9.1.5}\quad 设 $V$ 是 $n$ 维实列向量空间，$G$ 是 $n$ 阶正定实对称阵，对 $\alpha,\beta\in V$，定义
\[
  (\alpha,\beta)=\alpha'G\beta,
\]
我们来证明 $V$ 在上式的定义下成为欧氏空间. 定义 9.1.1 中的 (2),(3) 显然成立. 对 (1)，注意到 $\alpha'G\beta$ 是实数，其转置仍是它自己，而 $G$ 是对称阵，故
\[
  (\alpha,\beta)=\alpha'G\beta=(\alpha'G\beta)'=\beta'G'\alpha=\beta'G\alpha=(\beta,\alpha).
\]
又从 $G$ 是正定阵即可知道 (4) 成立.

当 $G=I_n$ 为单位阵时，$V$ 上内积就是例 9.1.1 中的标准内积. 对实列向量空间，标准内积可用矩阵乘法表示为
\[
  (\alpha,\beta)=\alpha'\beta.
\]
对实行向量空间，标准内积也可表示为
\[
  (\alpha,\beta)=\alpha\beta'.
\]
对 $n$ 维复列向量空间 $U$，若有正定 Hermite 矩阵 $H$，也可定义 $U$ 上内积：
\[
  (\alpha,\beta)=\alpha'\overline H\overline\beta.
\]
当 $H=I_n$ 为单位阵时，$U$ 上内积就是例 9.1.2 中的标准内积. 对复列向量空间，标准内积可用矩阵乘法表示为
\[
  (\alpha,\beta)=\alpha'\overline\beta.
\]
对复行向量空间，标准内积也可表示为
\[
  (\alpha,\beta)=\overline\alpha\beta'.
\]

有了内积的概念，我们便可定义向量的长度（或称范数）了.

\noindent\textbf{定义 9.1.3}\quad 设 $V$ 是实或复的内积空间，$\alpha$ 是 $V$ 中的向量，定义 $\alpha$ 的长度（或范数）为
\[
  \|\alpha\|=(\alpha,\alpha)^{\frac12},
\]
即实数 $(\alpha,\alpha)$ 的算术平方根.

注意由规则 (4) 可知，$(\alpha,\alpha)$ 总是非负实数. 从长度的定义知，$\|\alpha\|=0$ 当且仅当 $\alpha=0$. 当 $V=\mathbb R^n$ 且内积为标准内积时，若 $\alpha=(x_1,x_2,\cdots,x_n)$，则
\begin{equation}
  \|\alpha\|=\sqrt{x_1^2+x_2^2+\cdots+x_n^2}.
  \tag{9.1.2}
\end{equation}
利用范数可定义内积空间中两个向量的距离. 设 $\alpha,\beta\in V$，定义 $\alpha$ 与 $\beta$ 的距离为
\[
  d(\alpha,\beta)=\|\alpha-\beta\|.
\]
显然 $d(\alpha,\beta)=d(\beta,\alpha)$.

\noindent\textbf{定理 9.1.1}\quad 设 $V$ 是实或复的内积空间，$\alpha,\beta\in V$，$c$ 是任一常数（实数或复数），则
\[
\begin{gathered}
  (1)\quad \|c\alpha\|=|c|\|\alpha\|;\\
  (2)\quad |(\alpha,\beta)|\le \|\alpha\|\cdot\|\beta\|;\\
  (3)\quad \|\alpha+\beta\|\le \|\alpha\|+\|\beta\|.
\end{gathered}
\]

\noindent\textbf{证明}\quad (1) $\|c\alpha\|^2=(c\alpha,c\alpha)=c\overline c(\alpha,\alpha)=|c|^2\|\alpha\|^2$，故 $\|c\alpha\|=|c|\|\alpha\|$.

(2) 若 $\alpha=0$，则 $(0,\beta)=(0+0,\beta)=2(0,\beta)$，故 $(0,\beta)=0$，因此 (2) 成立.

若 $\alpha\ne0$，令
\[
  v=\beta-\frac{(\beta,\alpha)}{\|\alpha\|^2}\alpha,
\]
则 $(v,\alpha)=0$，且
\[
\begin{aligned}
  0\le \|v\|^2
  &=
  \left(\beta-\frac{(\beta,\alpha)}{\|\alpha\|^2}\alpha,
  \beta-\frac{(\beta,\alpha)}{\|\alpha\|^2}\alpha\right)\\
  &=(\beta,\beta)-\frac{(\beta,\alpha)}{\|\alpha\|^2}(\alpha,\beta)\\
  &=\|\beta\|^2-\frac{|(\alpha,\beta)|^2}{\|\alpha\|^2},
\end{aligned}
\]
由此即可得 (2).

(3) 我们有
\[
\begin{aligned}
  \|\alpha+\beta\|^2
  &=(\alpha+\beta,\alpha+\beta)\\
  &=\|\alpha\|^2+(\alpha,\beta)+(\beta,\alpha)+\|\beta\|^2\\
  &=\|\alpha\|^2+\|\beta\|^2+(\alpha,\beta)+\overline{(\alpha,\beta)}.
\end{aligned}
\]
由 (2) 得
\[
  |(\alpha,\beta)|\le \|\alpha\|\|\beta\|,
  \quad
  |\overline{(\alpha,\beta)}|\le \|\alpha\|\|\beta\|,
\]
故
\[
  \|\alpha+\beta\|^2\le \|\alpha\|^2+\|\beta\|^2+2\|\alpha\|\|\beta\|=(\|\alpha\|+\|\beta\|)^2. \quad \square
\]

\noindent\textbf{定义 9.1.4}\quad 当 $V$ 是实内积空间时，定义非零向量 $\alpha,\beta$ 的夹角 $\theta$ 之余弦为
\begin{equation}
  \cos\theta=\frac{(\alpha,\beta)}{\|\alpha\|\|\beta\|}.
  \tag{9.1.3}
\end{equation}
当 $V$ 是复内积空间时，定义非零向量 $\alpha,\beta$ 的夹角 $\theta$ 之余弦为
\[
  \cos\theta=\frac{|(\alpha,\beta)|}{\|\alpha\|\|\beta\|}.
\]
内积空间中两个向量 $\alpha,\beta$ 若适合 $(\alpha,\beta)=0$，则称 $\alpha$ 与 $\beta$ 垂直或正交，我们用记号 $\alpha\perp\beta$ 来表示. 显然，零向量和任何向量都正交；若 $\alpha$ 与 $\beta$ 正交，则 $\beta$ 也与 $\alpha$ 正交；两个非零向量 $\alpha,\beta$ 正交时夹角为 $90^\circ$.

定理 9.1.1 中的 (2) 通常称为 Cauchy-Schwarz 不等式，(3) 通常称为三角不等式. 读者注意在 (9.1.3) 式中要使 $\theta$ 有意义，必须保证 $|\cos\theta|\le1$，而这就是定理 9.1.1 中的 (2). 在定理 9.1.1 (3) 的证明中我们可看出：若 $\alpha$ 与 $\beta$ 正交，则 $(\alpha,\beta)=(\beta,\alpha)=0$，因此
\begin{equation}
  \|\alpha+\beta\|^2=\|\alpha\|^2+\|\beta\|^2.
  \tag{9.1.4}
\end{equation}
上式通常称为勾股定理，它是平面几何中勾股定理的推广.

设 $V$ 是 $n$ 维实行向量空间，内积取标准内积，从定理 9.1.1 的 (2) 立即可得到下列 Cauchy 不等式：
\[
  (x_1y_1+x_2y_2+\cdots+x_ny_n)^2
  \le
  (x_1^2+x_2^2+\cdots+x_n^2)(y_1^2+y_2^2+\cdots+y_n^2).
\]
设 $V$ 是由 $[a,b]$ 区间上连续函数全体构成的实线性空间，内积如例 9.1.4，则从定理 9.1.1 的 (2) 可得下列 Schwarz 不等式：
\[
  \left(\int_a^b f(t)g(t)\mathrm dt\right)^2
  \le
  \int_a^b f(t)^2\mathrm dt\int_a^b g(t)^2\mathrm dt.
\]

\section*{习题 9.1}

1. 设 $V$ 是内积空间，$\alpha\in V$，若对任意的 $\beta\in V$ 总有 $(\alpha,\beta)=0$，求证：$\alpha=0$.

2. 设 $V$ 是 $n$ 维内积空间，$\{e_1,e_2,\cdots,e_n\}$ 是 $V$ 的一组基，若
\[
  (\alpha,e_i)=(\beta,e_i),\quad i=1,2,\cdots,n,
\]
求证：$\alpha=\beta$.

3. 若 $(\ ,\ )_1,(\ ,\ )_2$ 是 $V$ 上的两个内积，证明：
\[
  (\alpha,\beta)=(\alpha,\beta)_1+(\alpha,\beta)_2
\]
定义了 $V$ 上的另一个内积.

4. 取 $n$ 维行向量空间 $\mathbb R^n$ 上的标准内积，标准单位行向量 $\{e_1,e_2,\cdots,e_n\}$ 构成 $\mathbb R^n$ 的一组基. 求证：对任意的 $\alpha\in\mathbb R^n$，有
\[
  \alpha=(\alpha,e_1)e_1+(\alpha,e_2)e_2+\cdots+(\alpha,e_n)e_n.
\]

5. 证明：在 $\mathbb R^2$（取标准内积）中存在一个非零线性变换 $\varphi$，使 $\varphi(\alpha)\perp\alpha$ 对任一 $\alpha\in\mathbb R^2$ 成立；但在 $\mathbb C^2$（取标准内积）中这样的非零线性变换不存在.

6. 证明：在内积空间中平行四边形两对角线平方和等于四边平方和，即
\[
  \|\alpha+\beta\|^2+\|\alpha-\beta\|^2=2\|\alpha\|^2+2\|\beta\|^2.
\]

7. 设 $V$ 是实系数多项式全体构成的实线性空间，对任意的实系数多项式
\[
  f(x)=a_0+a_1x+\cdots+a_nx^n,\quad
  g(x)=b_0+b_1x+\cdots+b_mx^m,
\]
定义
\[
  (f,g)=\sum_{i,j}\frac{a_ib_j}{i+j+1},
\]
证明：$(f,g)$ 定义了 $V$ 上的内积.

8. 证明：下列 $(A,B)$ 分别定义了 $n$ 阶实矩阵空间和 $n$ 阶复矩阵空间上的内积（称为 Frobenius 内积）：

(1) 设 $V$ 为 $n$ 阶实矩阵空间，对任意的 $n$ 阶实矩阵 $A=(a_{ij}),B=(b_{ij})$，定义
\[
  (A,B)=\operatorname{tr}(AB')=\sum_{i=1}^n\sum_{j=1}^n a_{ij}b_{ij};
\]

(2) 设 $V$ 为 $n$ 阶复矩阵空间，对任意的 $n$ 阶复矩阵 $A=(a_{ij}),B=(b_{ij})$，定义
\[
  (A,B)=\operatorname{tr}(A\overline B')=\sum_{i=1}^n\sum_{j=1}^n a_{ij}\overline{b_{ij}}.
\]

\section*{§9.2\quad 内积的表示和正交基}

这一节我们将只考虑有限维内积空间. 我们要讨论的第一个问题是：给定内积空间的一组基以后，如何用坐标向量来表示向量的内积. 具体来说，设 $V$ 是欧氏空间（酉空间），$\{v_1,v_2,\cdots,v_n\}$ 是 $V$ 的一组基. 如果 $(v_i,v_j)=g_{ij}(i,j=1,2,\cdots,n)$，$\alpha=a_1v_1+a_2v_2+\cdots+a_nv_n$，$\beta=b_1v_1+b_2v_2+\cdots+b_nv_n$，问：$(\alpha,\beta)$ 等于什么？

用向量内积的性质，很容易给出这个问题的答案. 当 $V$ 是欧氏空间时，
\[
  (\alpha,\beta)=\left(\sum_{i=1}^n a_iv_i,\sum_{j=1}^n b_jv_j\right)
  =\sum_{i,j=1}^n a_ig_{ij}b_j.
\]
我们把上述结论写成矩阵形式：
\[
  (\alpha,\beta)=
  (a_1,a_2,\cdots,a_n)
  \left(\begin{array}{cccc}
    g_{11} & g_{12} & \cdots & g_{1n}\\
    g_{21} & g_{22} & \cdots & g_{2n}\\
    \vdots & \vdots & & \vdots\\
    g_{n1} & g_{n2} & \cdots & g_{nn}
  \end{array}\right)
  \left(\begin{array}{c}
    b_1\\ b_2\\ \vdots\\ b_n
  \end{array}\right),
\]
其中矩阵
\[
  G=
  \left(\begin{array}{cccc}
    g_{11} & g_{12} & \cdots & g_{1n}\\
    g_{21} & g_{22} & \cdots & g_{2n}\\
    \vdots & \vdots & & \vdots\\
    g_{n1} & g_{n2} & \cdots & g_{nn}
  \end{array}\right)
  =
  \left(\begin{array}{cccc}
    (v_1,v_1) & (v_1,v_2) & \cdots & (v_1,v_n)\\
    (v_2,v_1) & (v_2,v_2) & \cdots & (v_2,v_n)\\
    \vdots & \vdots & & \vdots\\
    (v_n,v_1) & (v_n,v_2) & \cdots & (v_n,v_n)
  \end{array}\right)
\]
称为基向量 $\{v_1,v_2,\cdots,v_n\}$ 的 Gram（格列姆）矩阵或内积空间 $V$ 在给定基下的度量矩阵.

于是，我们得到了内积在给定基下的表示：
\begin{equation}
  (\alpha,\beta)=x'Gy,
  \tag{9.2.1}
\end{equation}
其中 $x,y$ 分别是向量 $\alpha,\beta$ 在给定基下的坐标向量.

再来看矩阵 $G$. 因为 $(v_i,v_j)=(v_j,v_i)$，所以 $G$ 是实对称阵. 又因为对任意的非零向量 $\alpha$，总有 $(\alpha,\alpha)>0$，所以 $x'Gx>0$ 对一切 $n$ 维非零实列向量 $x$ 成立. 这表明 $G$ 是一个正定阵. 反之，若给定 $n$ 阶正定实对称阵 $G$，利用 (9.2.1) 式也可以定义 $V$ 上的内积（参考例 9.1.5）. 由此我们可以看出，若给定了 $n$ 维实线性空间 $V$ 的一组基，则 $V$ 上的内积结构和 $n$ 阶正定实对称阵之间存在着一个一一对应.

设 $V$ 是酉空间，我们类似可证若 $\alpha=a_1v_1+a_2v_2+\cdots+a_nv_n,\beta=b_1v_1+b_2v_2+\cdots+b_nv_n$，令
\[
  H=
  \left(\begin{array}{cccc}
    (v_1,v_1) & (v_1,v_2) & \cdots & (v_1,v_n)\\
    (v_2,v_1) & (v_2,v_2) & \cdots & (v_2,v_n)\\
    \vdots & \vdots & & \vdots\\
    (v_n,v_1) & (v_n,v_2) & \cdots & (v_n,v_n)
  \end{array}\right),
\]
则
\begin{equation}
  (\alpha,\beta)=x'\overline H\overline y,
  \tag{9.2.2}
\end{equation}
其中 $x,y$ 分别是 $\alpha,\beta$ 的坐标向量，$H$ 是一个正定 Hermite 矩阵.

从 (9.2.1) 式和 (9.2.2) 式，我们自然地会提出这样一个问题：在 $V$ 中是否存在这样一组基，它的 Gram 矩阵是单位阵 $I_n$？如果存在，那么在这组基下的内积表示将特别简单. 首先，我们来看如果 $G=I_n$（或 $H=I_n$），基向量要满足什么条件？显然，$G$（或 $H$）等于 $I_n$ 的充分必要条件是：
\[
  (v_i,v_j)=0\ (i\ne j),\quad (v_i,v_i)=1.
\]

\noindent\textbf{定义 9.2.1}\quad 设 $\{e_1,e_2,\cdots,e_n\}$ 是 $n$ 维内积空间 $V$ 的一组基. 若 $e_i\perp e_j$ 对一切 $i\ne j$ 成立，则称这组基是 $V$ 的一组正交基. 又若 $V$ 的一组正交基中每个基向量的长度都等于 1，则称这组正交基为标准正交基.

显然在标准正交基下，度量矩阵就是单位阵，那么在任意的有限维内积空间中，标准正交基一定存在吗？容易验证，$n$ 维行（列）向量空间在标准内积下，它的标准基（即由标准单位行（列）向量组成的基）是标准正交基. 在证明一般内积空间的标准正交基的存在性之前，我们先证明正交向量组的两个重要性质.

\noindent\textbf{引理 9.2.1}\quad 内积空间 $V$ 中的任意一组两两正交的非零向量必线性无关.

\noindent\textbf{证明}\quad 设 $v_1,v_2,\cdots,v_m$ 是 $V$ 中两两正交的非零向量，若
\[
  k_1v_1+k_2v_2+\cdots+k_mv_m=0,
\]
则对任一 $1\le i\le m$，有
\[
  (k_1v_1+k_2v_2+\cdots+k_mv_m,v_i)=0.
\]
由于 $v_i\perp v_j(i\ne j)$，故由上式可得 $k_i(v_i,v_i)=0$，又 $v_i\ne0$，从而 $k_i=0$. $\square$

\noindent\textbf{引理 9.2.2}\quad 设向量 $\alpha$ 和 $\beta_1,\beta_2,\cdots,\beta_k$ 都正交，则 $\alpha$ 和 $L(\beta_1,\beta_2,\cdots,\beta_k)$ 中的每个向量都正交.

\noindent\textbf{证明}\quad 任取 $\beta=b_1\beta_1+b_2\beta_2+\cdots+b_k\beta_k\in L(\beta_1,\beta_2,\cdots,\beta_k)$，则
\[
  (\beta,\alpha)=(b_1\beta_1+b_2\beta_2+\cdots+b_k\beta_k,\alpha)
  =\sum_{i=1}^k b_i(\beta_i,\alpha)=0,
\]
结论得证. $\square$

\noindent\textbf{推论 9.2.1}\quad $n$ 维内积空间中任意一个正交非零向量组的向量个数不超过 $n$.

若 $\alpha$ 是一个非零向量，则 $\dfrac1{\|\alpha\|}\alpha$ 是一个长度等于 1 的向量，这个过程称为单位化. 因此要证明存在标准正交基，我们只需先证明存在正交基，然后将正交基单位化即可. 我们从一组线性无关的向量出发，设法构造出一组正交向量来. 设 $u_1,u_2,\cdots,u_m$ 是 $V$ 中 $m$ 个线性无关的向量，我们采用数学归纳法. 当只有一个向量时，令 $v_1=u_1$. 假设 $v_1,v_2,\cdots,v_k$ 已经构造好，是一组两两正交的非零向量，我们用待定系数法求 $v_{k+1}$. 设
\[
  v_{k+1}=u_{k+1}+a_1v_1+a_2v_2+\cdots+a_kv_k,
\]
两边和 $v_j(j\le k)$ 作内积，便得到
\[
  0=(v_{k+1},v_j)=(u_{k+1},v_j)+a_j(v_j,v_j),
\]
于是
\[
  a_j=-\frac{(u_{k+1},v_j)}{(v_j,v_j)}.
\]
下面我们给出定理及其证明.

\noindent\textbf{定理 9.2.1}\quad 设 $V$ 是内积空间，$u_1,u_2,\cdots,u_m$ 是 $V$ 中 $m$ 个线性无关的向量，则在 $V$ 中存在 $m$ 个两两正交的非零向量 $v_1,v_2,\cdots,v_m$，使由 $v_1,v_2,\cdots,v_m$ 张成的子空间恰好为由 $u_1,u_2,\cdots,u_m$ 张成的子空间，即 $v_1,v_2,\cdots,v_m$ 是该子空间的一组正交基.

\noindent\textbf{证明}\quad 设 $v_1=u_1$，其余 $v_i$ 可用数学归纳法定义如下：假设 $v_1,\cdots,v_k(k<m)$ 已定义好，这时 $v_1,\cdots,v_k$ 两两正交非零且 $L(v_1,\cdots,v_k)=L(u_1,\cdots,u_k)$. 令
\begin{equation}
  v_{k+1}=u_{k+1}-\sum_{j=1}^k\frac{(u_{k+1},v_j)}{\|v_j\|^2}v_j.
  \tag{9.2.3}
\end{equation}
注意 $v_{k+1}\ne0$，否则 $u_{k+1}$ 将是 $v_1,\cdots,v_k$ 的线性组合，从而也是 $u_1,\cdots,u_k$ 的线性组合，此与 $u_1,u_2,\cdots,u_m$ 线性无关矛盾. 又对任意的 $1\le i\le k$，有
\[
\begin{aligned}
  (v_{k+1},v_i)
  &=(u_{k+1},v_i)-\sum_{j=1}^k\frac{(u_{k+1},v_j)}{\|v_j\|^2}(v_j,v_i)\\
  &=(u_{k+1},v_i)-(u_{k+1},v_i)=0,
\end{aligned}
\]
因此 $v_1,\cdots,v_k,v_{k+1}$ 两两正交. 由 (9.2.3) 式可知 $u_{k+1}\in L(v_1,\cdots,v_k,v_{k+1})$ 及
\[
  v_{k+1}\in L(v_1,\cdots,v_k)+L(u_{k+1})
  =L(u_1,\cdots,u_k)+L(u_{k+1})=L(u_1,\cdots,u_k,u_{k+1}),
\]
于是 $L(v_1,\cdots,v_k,v_{k+1})=L(u_1,\cdots,u_k,u_{k+1})$，这就证明了结论. $\square$

上述正交化过程通常称为 Gram-Schmidt（格列姆--施密特）方法.

\noindent\textbf{推论 9.2.2}\quad 任一有限维内积空间均有标准正交基.

\noindent\textbf{例 9.2.1}\quad 设 $V$ 是三维实行向量空间，内积为标准内积. 又已知 3 个线性无关的向量：
\[
  u_1=(3,0,4),\quad u_2=(-1,0,7),\quad u_3=(2,9,11),
\]
用 Gram-Schmidt 方法求 $V$ 的一组标准正交基.

\noindent\textbf{解}\quad 令
\[
\begin{aligned}
  v_1&=(3,0,4),\\
  v_2&=(-1,0,7)-\frac{-1\cdot3+0+7\cdot4}{25}(3,0,4)=(-4,0,3),\\
  v_3&=(2,9,11)-\frac{2\cdot3+0+11\cdot4}{25}(3,0,4)\\
  &\quad-\frac{2\cdot(-4)+0+11\cdot3}{25}(-4,0,3)\\
  &=(0,9,0).
\end{aligned}
\]
再令
\[
  w_1=\frac15(3,0,4),\quad w_2=\frac15(-4,0,3),\quad w_3=(0,1,0),
\]
则 $\{w_1,w_2,w_3\}$ 是 $V$ 的一组标准正交基.

下面我们要讨论子空间的正交问题.

\noindent\textbf{定义 9.2.2}\quad 设 $U$ 是内积空间 $V$ 的子空间，令
\[
  U^\perp=\{v\in V\mid (v,U)=0\},
\]
这里 $(v,U)=0$ 表示对一切 $u\in U$，均有 $(v,u)=0$. 容易验证 $U^\perp$ 是 $V$ 的子空间，称为 $U$ 的正交补空间.

\noindent\textbf{定理 9.2.2}\quad 设 $V$ 是 $n$ 维内积空间，$U$ 是 $V$ 的子空间，则

(1) $V=U\oplus U^\perp$;

(2) $U$ 的任一组标准正交基均可扩张为 $V$ 的一组标准正交基.

\noindent\textbf{证明}\quad (1) 若 $x\in U\cap U^\perp$，则 $(x,x)=0$，因此 $x=0$，即 $U\cap U^\perp=0$. 另一方面，由推论 9.2.2 可知，存在 $U$ 的一组标准正交基 $\{e_1,e_2,\cdots,e_m\}$. 对任意的 $v\in V$，令
\[
  u=(v,e_1)e_1+(v,e_2)e_2+\cdots+(v,e_m)e_m,
\]
则 $u\in U$. 又令 $w=v-u$，则对任一 $e_i(i=1,2,\cdots,m)$，有
\[
  (w,e_i)=(v,e_i)-(u,e_i)=(v,e_i)-(v,e_i)=0.
\]
因此 $w\in U^\perp$，又 $v=u+w$，这就证明了 $V=U\oplus U^\perp$.

(2) 设 $\{e_1,e_2,\cdots,e_m\}$ 是 $U$ 的任一组标准正交基，$\{e_{m+1},\cdots,e_n\}$ 是 $U^\perp$ 的任一组标准正交基，则显然 $\{e_1,e_2,\cdots,e_n\}$ 是 $V$ 的一组标准正交基. $\square$

\noindent\textbf{定义 9.2.3}\quad 设 $V$ 是 $n$ 维内积空间，$V_1,V_2,\cdots,V_k$ 是 $V$ 的子空间. 如果对任意的 $\alpha\in V_i$ 和任意的 $\beta\in V_j$，均有 $(\alpha,\beta)=0$，则称子空间 $V_i$ 和 $V_j$ 正交. 若 $V=V_1+V_2+\cdots+V_k$ 且 $V_i$ 两两正交，则称 $V$ 是 $V_1,V_2,\cdots,V_k$ 的正交和，记为
\[
  V=V_1\perp V_2\perp\cdots\perp V_k.
\]

\noindent\textbf{引理 9.2.3}\quad 正交和必为直和且任一 $V_i$ 和其余子空间的和正交.

\noindent\textbf{证明}\quad 对任意的 $v_i\in V_i$ 和 $\sum_{j\ne i}v_j(v_j\in V_j)$，有
\[
  \left(v_i,\sum_{j\ne i}v_j\right)=\sum_{j\ne i}(v_i,v_j)=0,
\]
因此后一个结论成立. 任取 $v\in V_i\cap\left(\sum_{j\ne i}V_j\right)$，则由上述论证可得 $(v,v)=0$，故 $v=0$，从而 $V_i\cap\left(\sum_{j\ne i}V_j\right)=0$，即正交和必为直和. $\square$

由于上述引理，正交和通常也称为正交直和.

\noindent\textbf{定义 9.2.4}\quad 设 $V=V_1\perp V_2\perp\cdots\perp V_k$，定义 $V$ 上的线性变换 $E_i(i=1,2,\cdots,k)$ 如下：若 $v=v_1+\cdots+v_i+\cdots+v_k(v_i\in V_i)$，令 $E_i(v)=v_i$. 容易验证 $E_i$ 是 $V$ 上的线性变换，且满足
\[
  E_i^2=E_i,\quad E_iE_j=0\ (i\ne j),\quad E_1+E_2+\cdots+E_k=I_V.
\]
线性变换 $E_i$ 称为 $V$ 到 $V_i$ 上的正交投影（简称投影）.

\noindent\textbf{命题 9.2.1}\quad 设 $U$ 是内积空间 $V$ 的子空间，$V=U\perp U^\perp$. 设 $E$ 是 $V$ 到 $U$ 上的正交投影，则对任意的 $\alpha,\beta\in V$，有
\[
  (E(\alpha),\beta)=(\alpha,E(\beta)).
\]

\noindent\textbf{证明}\quad 设 $\alpha=u_1+w_1,\beta=u_2+w_2$，其中 $u_1,u_2\in U,w_1,w_2\in U^\perp$，则 $E(\alpha)=u_1,E(\beta)=u_2$，于是
\[
  (E(\alpha),\beta)=(u_1,u_2+w_2)=(u_1,u_2)+(u_1,w_2)=(u_1,u_2),
\]
\[
  (\alpha,E(\beta))=(u_1+w_1,u_2)=(u_1,u_2)+(w_1,u_2)=(u_1,u_2).
\]
由此即得结论. $\square$

下面的结论是``斜边大于直角边''这一几何命题在内积空间中的推广.

\noindent\textbf{命题 9.2.2 (Bessel（贝塞尔）不等式)}\quad 设 $v_1,v_2,\cdots,v_m$ 是内积空间 $V$ 中的正交非零向量组，$y$ 是 $V$ 中任一向量，则
\[
  \sum_{k=1}^m\frac{|(y,v_k)|^2}{\|v_k\|^2}\le \|y\|^2,
\]
且等号成立的充分必要条件是 $y$ 属于由 $\{v_1,v_2,\cdots,v_m\}$ 张成的子空间.

\noindent\textbf{证明}\quad 令
\[
  x=\sum_{k=1}^m\frac{(y,v_k)}{\|v_k\|^2}v_k,
\]
则 $x$ 属于由 $\{v_1,v_2,\cdots,v_m\}$ 张成的子空间. 容易验证
\[
  (y-x,v_k)=0,\quad k=1,2,\cdots,m,
\]
因此 $(y-x,x)=0$. 由勾股定理可得
\[
  \|y\|^2=\|y-x\|^2+\|x\|^2,
\]
故
\[
  \|x\|^2\le \|y\|^2.
\]
又由 $v_1,v_2,\cdots,v_m$ 两两正交不难算出
\[
  \|x\|^2=\sum_{k=1}^m\frac{|(y,v_k)|^2}{\|v_k\|^2}.
\]
若 $y$ 属于由 $\{v_1,v_2,\cdots,v_m\}$ 张成的子空间，则 $y=x$，故等号成立. 反之，若等号成立，则 $\|y-x\|^2=0$，故 $y=x$，即 $y$ 属于由 $\{v_1,v_2,\cdots,v_m\}$ 张成的子空间. $\square$

Bessel 不等式在数学分析中有重要的应用.

\section*{习题 9.2}

1. 用 Gram-Schmidt 方法求由下列向量张成的子空间的标准正交基：

(1) $u_1=(1,1,1),u_2=(1,1,0),u_3=(1,0,0)$;

(2) $u_1=(1,2,2,-1),u_2=(1,1,-5,3),u_3=(3,2,8,-7)$;

(3) $u_1=(1,1,-1,-2),u_2=(5,8,-2,-3),u_3=(3,9,3,8)$.

2. 设 $e_1,e_2,\cdots,e_n$ 为内积空间 $V$ 的一组标准正交基，求证：对任意的 $\alpha,\beta\in V$，有
\[
  (\alpha,\beta)=(\alpha,e_1)\overline{(\beta,e_1)}
  +(\alpha,e_2)\overline{(\beta,e_2)}
  +\cdots
  +(\alpha,e_n)\overline{(\beta,e_n)}.
\]

3. 设 $V$ 是内积空间，$U$ 是 $V$ 的子空间，若 $U$ 由向量组 $\{u_1,u_2,\cdots,u_m\}$ 生成，求证：
\[
  U^\perp=\{v\in V\mid (v,u_i)=0,\ i=1,2,\cdots,m\}.
\]

4. 设 $U,U_1,U_2$ 是有限维内积空间 $V$ 的子空间，求证：
\[
\begin{gathered}
  (1)\quad (U^\perp)^\perp=U;\\
  (2)\quad (U_1+U_2)^\perp=U_1^\perp\cap U_2^\perp;\\
  (3)\quad (U_1\cap U_2)^\perp=U_1^\perp+U_2^\perp;\\
  (4)\quad V^\perp=0,\quad 0^\perp=V.
\end{gathered}
\]

5. 设 $S$ 是有限维内积空间 $V$ 的子集，证明：

(1) $S^\perp=\{\alpha\in V\mid (\alpha,S)=0\}$ 是 $V$ 的子空间；

(2) $(S^\perp)^\perp$ 等于由 $S$ 生成的子空间.

6. 设线性子空间 $U$ 为下列线性方程组的解空间：
\[
  \begin{cases}
    2x_1+x_2+3x_3-x_4=0,\\
    3x_1+2x_2-2x_4=0,\\
    3x_1+x_2+9x_3-x_4=0,
  \end{cases}
\]
试求 $U^\perp$ 适合的线性方程组.

7. 设 $V$ 为 $n$ 阶实矩阵空间（取 Frobenius 内积），$V_1$ 是 $n$ 阶实对称阵构成的子空间，$V_2$ 是 $n$ 阶实反对称阵构成的子空间，求证：
\[
  V=V_1\perp V_2.
\]

8. 设 $V$ 是 $n$ 维欧氏空间，$\{e_1,e_2,\cdots,e_n\}$ 是 $V$ 的一组基，$c_1,c_2,\cdots,c_n$ 是 $n$ 个实数，求证：存在唯一的向量 $\alpha\in V$，使
\[
  (\alpha,e_i)=c_i,\quad i=1,2,\cdots,n.
\]

9. 设 $u_1,u_2,\cdots,u_m$ 是欧氏空间 $V$ 中的 $m$ 个向量，矩阵
\[
  G(u_1,u_2,\cdots,u_m)=
  \left(\begin{array}{cccc}
    (u_1,u_1) & (u_1,u_2) & \cdots & (u_1,u_m)\\
    (u_2,u_1) & (u_2,u_2) & \cdots & (u_2,u_m)\\
    \vdots & \vdots & & \vdots\\
    (u_m,u_1) & (u_m,u_2) & \cdots & (u_m,u_m)
  \end{array}\right)
\]
称为 $u_1,u_2,\cdots,u_m$ 的 Gram 矩阵. 证明：若用 Gram-Schmidt 方法将线性无关向量组 $\{u_1,u_2,\cdots,u_m\}$ 变成正交向量组 $\{v_1,v_2,\cdots,v_m\}$，则这两组向量的 Gram 矩阵的行列式值不变，即
\[
  |G(u_1,u_2,\cdots,u_m)|=|G(v_1,v_2,\cdots,v_m)|
  =\|v_1\|^2\|v_2\|^2\cdots\|v_m\|^2.
\]

10. 证明：向量 $u_1,u_2,\cdots,u_m$ 线性无关的充分必要条件是它们的 Gram 矩阵为非异阵.

11. 证明下列不等式：
\[
  0\le |G(u_1,u_2,\cdots,u_m)|\le \|u_1\|^2\|u_2\|^2\cdots\|u_m\|^2,
\]
后一个等号成立的充分必要条件是 $u_i$ 两两正交或者某个 $u_i=0$.

12. 设 $A=(a_{ij})$ 是 $n$ 阶实矩阵，证明下列 Hadamard（阿达玛）不等式：
\[
  |A|^2\le \prod_{j=1}^n\sum_{i=1}^n a_{ij}^2.
\]

13. 令 $V$ 是区间 $[-1,1]$ 上次数不超过 $n$ 的实系数多项式构成的实线性空间，$V$ 上的内积定义为
\[
  (f,g)=\int_{-1}^1 f(x)g(x)\mathrm dx,
\]
证明：$V$ 是 $n+1$ 维欧氏空间，并且下列多项式组成 $V$ 的一组正交基：
\[
  P_0(x)=1,\quad
  P_k(x)=\frac1{2^kk!}\cdot\frac{\mathrm d^k}{\mathrm dx^k}\left[(x^2-1)^k\right]\quad(k=1,2,\cdots,n),
\]
称之为 Legendre（勒让德）多项式.

\section*{§9.3\quad 伴\quad 随}

在这一节里我们将考察内积空间中的线性变换，内积空间中的线性变换常常被称为线性算子.

现设 $V$ 是 $n$ 维内积空间（不妨设之为酉空间），取 $V$ 的一组标准正交基 $\{e_1,e_2,\cdots,e_n\}$. 假设 $\varphi$ 是 $V$ 上的线性变换，且它在这组基下的表示矩阵是
\[
  A=
  \left(\begin{array}{cccc}
    a_{11} & a_{12} & \cdots & a_{1n}\\
    a_{21} & a_{22} & \cdots & a_{2n}\\
    \vdots & \vdots & & \vdots\\
    a_{n1} & a_{n2} & \cdots & a_{nn}
  \end{array}\right).
\]
又设 $\alpha=a_1e_1+a_2e_2+\cdots+a_ne_n,\beta=b_1e_1+b_2e_2+\cdots+b_ne_n$. 记 $x=(a_1,a_2,\cdots,a_n)',y=(b_1,b_2,\cdots,b_n)'$ 分别是 $\alpha,\beta$ 的坐标向量，则
\begin{equation}
  (\varphi(\alpha),\beta)=(Ax)'\overline y=x'A'\overline y=x'(\overline A'y).
  \tag{9.3.1}
\end{equation}
记矩阵 $\overline A'$ 在 $V$ 上定义的线性变换为 $\psi$，即对任意的 $\beta=b_1e_1+b_2e_2+\cdots+b_ne_n$，定义 $\psi(\beta)=c_1e_1+c_2e_2+\cdots+c_ne_n$，其中
\[
  c_i=\overline{a_{1i}}b_1+\overline{a_{2i}}b_2+\cdots+\overline{a_{ni}}b_n\quad(i=1,2,\cdots,n),
\]
即 $\psi(\beta)$ 的坐标向量
\[
  \left(\begin{array}{c}
    c_1\\ c_2\\ \vdots\\ c_n
  \end{array}\right)
  =\overline A'y
  =
  \left(\begin{array}{cccc}
    \overline{a_{11}} & \overline{a_{21}} & \cdots & \overline{a_{n1}}\\
    \overline{a_{12}} & \overline{a_{22}} & \cdots & \overline{a_{n2}}\\
    \vdots & \vdots & & \vdots\\
    \overline{a_{1n}} & \overline{a_{2n}} & \cdots & \overline{a_{nn}}
  \end{array}\right)
  \left(\begin{array}{c}
    b_1\\ b_2\\ \vdots\\ b_n
  \end{array}\right).
\]
于是 (9.3.1) 式表明
\[
  (\varphi(\alpha),\beta)=(\alpha,\psi(\beta))
\]
对一切 $\alpha,\beta\in V$ 成立.

读者很容易验证上述结论对欧氏空间也成立.

\noindent\textbf{定义 9.3.1}\quad 设 $\varphi$ 是内积空间 $V$ 上的线性算子，若存在 $V$ 上的线性算子 $\varphi^*$，使等式
\[
  (\varphi(\alpha),\beta)=(\alpha,\varphi^*(\beta))
\]
对一切 $\alpha,\beta\in V$ 成立，则称 $\varphi^*$ 是 $\varphi$ 的伴随算子，简称为 $\varphi$ 的伴随.

\noindent\textbf{定理 9.3.1}\quad 设 $V$ 是 $n$ 维内积空间，$\varphi$ 是 $V$ 上的线性变换，则存在 $V$ 上唯一的线性变换 $\varphi^*$，使对一切 $\alpha,\beta\in V$，成立
\begin{equation}
  (\varphi(\alpha),\beta)=(\alpha,\varphi^*(\beta)).
  \tag{9.3.2}
\end{equation}

\noindent\textbf{证明}\quad 只需证明唯一性. 若 $\varphi^\#$ 是 $V$ 上的线性变换且
\[
  (\varphi(\alpha),\beta)=(\alpha,\varphi^\#(\beta))
\]
对一切 $\alpha,\beta\in V$ 成立，则 $(\alpha,\varphi^\#(\beta))=(\alpha,\varphi^*(\beta))$ 对一切 $\alpha\in V$ 成立，即
\[
  (\alpha,\varphi^\#(\beta)-\varphi^*(\beta))=0
\]
对一切 $\alpha\in V$ 成立，特别，对 $\alpha=\varphi^\#(\beta)-\varphi^*(\beta)$ 也成立. 由内积定义即知 $\varphi^\#(\beta)-\varphi^*(\beta)=0$，即 $\varphi^\#(\beta)=\varphi^*(\beta)$，而 $\beta$ 是任意的，故有 $\varphi^\#=\varphi^*$. $\square$

定理 9.3.1 表明，对有限维内积空间 $V$ 上的任一线性算子，它的伴随必存在且唯一. 从本节开始的分析我们还知道，线性算子和它的伴随在同一组标准正交基下的表示矩阵有下面的关系.

\noindent\textbf{定理 9.3.2}\quad 设 $V$ 是 $n$ 维内积空间，$\{e_1,e_2,\cdots,e_n\}$ 是 $V$ 的一组标准正交基. 若 $V$ 上的线性算子 $\varphi$ 在这组基下的表示矩阵为 $A$，则当 $V$ 是酉空间时，$\varphi^*$ 在同一组基下的表示矩阵为 $\overline A'$，即 $A$ 的共轭转置；当 $V$ 是欧氏空间时，$\varphi^*$ 的表示矩阵为 $A'$，即 $A$ 的转置.

\noindent\textbf{证明}\quad 由伴随的唯一性知道本节一开始由 $\overline A'$ 定义的线性变换 $\psi$ 就是 $\varphi$ 的伴随，而 $\psi$ 的表示矩阵就是 $\overline A'$. $\square$

伴随算子有下列性质.

\noindent\textbf{定理 9.3.3}\quad 设 $V$ 是有限维内积空间，若 $\varphi$ 及 $\psi$ 是 $V$ 上的线性变换，$c$ 为常数，则
\[
\begin{gathered}
  (1)\quad (\varphi+\psi)^*=\varphi^*+\psi^*;\\
  (2)\quad (c\varphi)^*=\overline c\varphi^*;\\
  (3)\quad (\varphi\psi)^*=\psi^*\varphi^*;\\
  (4)\quad (\varphi^*)^*=\varphi.
\end{gathered}
\]

\noindent\textbf{证明}\quad 由矩阵和线性变换的一一对应关系及矩阵共轭转置的性质即得. $\square$

关于伴随的不变子空间和特征值有下面的结果.

\noindent\textbf{命题 9.3.1}\quad 设 $V$ 是 $n$ 维内积空间，$\varphi$ 是 $V$ 上的线性算子.

(1) 若 $U$ 是 $\varphi$ 的不变子空间，则 $U^\perp$ 是 $\varphi^*$ 的不变子空间；

(2) 若 $\varphi$ 的全体特征值为 $\lambda_1,\lambda_2,\cdots,\lambda_n$，则 $\varphi^*$ 的全体特征值为 $\overline{\lambda_1},\overline{\lambda_2},\cdots,\overline{\lambda_n}$.

\noindent\textbf{证明}\quad (1) 任取 $\alpha\in U,\beta\in U^\perp$，因为
\[
  (\alpha,\varphi^*(\beta))=(\varphi(\alpha),\beta)=0,
\]
所以 $U^\perp$ 是 $\varphi^*$ 的不变子空间.

(2) 取 $V$ 的一组标准正交基，设 $\varphi$ 在这组基下的表示矩阵为 $A$，则无论 $V$ 是酉空间还是欧氏空间，$\varphi^*$ 的表示矩阵总可写为 $\overline A'$. 由假设
\[
  |\lambda I_n-A|=(\lambda-\lambda_1)(\lambda-\lambda_2)\cdots(\lambda-\lambda_n),
\]
则容易验证
\[
  |\lambda I_n-\overline A'|=(\lambda-\overline{\lambda_1})(\lambda-\overline{\lambda_2})\cdots(\lambda-\overline{\lambda_n}),
\]
故结论成立. $\square$

\section*{习题 9.3}

1. 设 $V$ 是二维复行向量空间（取标准内积），$\{e_1,e_2\}$ 是 $V$ 的一组标准正交基. 设 $V$ 上的线性算子 $\varphi$ 满足：
\[
  \varphi(e_1)=(1,-2),\quad \varphi(e_2)=(i,-1),
\]
求 $\varphi$ 及 $\varphi^*$ 在这组基下的表示矩阵. 又若 $x=(x_1,x_2)$，求 $\varphi^*(x)$.

2. 设 $V$ 是有限维内积空间，$\varphi$ 是 $V$ 上的线性算子，求证：$\operatorname{Im}\varphi^*=(\operatorname{Ker}\varphi)^\perp$.

3. 设 $\varphi$ 是有限维内积空间 $V$ 上的线性算子，若 $\varphi$ 是线性自同构，求证：$\varphi^*$ 也是线性自同构且 $(\varphi^*)^{-1}=(\varphi^{-1})^*$.

4. 设 $\alpha,\beta$ 是 $n$ 维酉空间 $V$ 中的两个向量，$V$ 上的变换 $\varphi$ 定义为 $\varphi(x)=(x,\alpha)\beta$，求证：$\varphi$ 是 $V$ 上的线性变换，并求 $\varphi^*$. 若 $\alpha,\beta$ 是两个正交单位向量，将它们扩展为 $V$ 的一组标准正交基 $\{e_1=\alpha,e_2=\beta,e_3,\cdots,e_n\}$，求 $\varphi$ 和 $\varphi^*$ 在这组基下的表示矩阵.

5. 设 $V$ 是由 $n$ 阶实矩阵全体构成的欧氏空间（取 Frobenius 内积），$T$ 是一个 $n$ 阶实矩阵，定义 $V$ 上的线性变换 $\varphi(A)=TA$，求 $\varphi$ 的伴随.

6. 设向量 $x$ 既是线性算子 $\varphi$ 的属于特征值 $\lambda_1$ 的特征向量，又是 $\varphi^*$ 的属于特征值 $\lambda_2$ 的特征向量，证明：$\lambda_1=\overline{\lambda_2}$.

7. 设 $\varphi$ 是酉空间 $V$ 上的线性算子，$\varphi$ 的极小多项式为 $g(x)$，证明：$\varphi^*$ 的极小多项式为 $\overline{g}(x)$，这里 $\overline{g}(x)$ 的系数等于 $g(x)$ 系数的共轭.

8. 设 $\varphi$ 是有限维内积空间 $V$ 上的线性算子，$U$ 是 $V$ 的子空间. 若 $U$ 既是 $\varphi$ 的不变子空间，又是 $\varphi^*$ 的不变子空间，求证：$(\varphi|_U)^*=\varphi^*|_U$.

\section*{§9.4\quad 内积空间的同构、正交变换和酉变换}

在第四章中，我们曾经讨论过抽象的 $n$ 维线性空间和 $n$ 维列向量空间（或行向量空间）的同构. 我们知道若在 $V$ 中取定一组基，将 $V$ 中向量映射到它在这组基下坐标向量的映射是一个线性同构. 现在设 $V$ 是 $n$ 维欧氏空间，$\{v_1,v_2,\cdots,v_n\}$ 是 $V$ 的一组标准正交基. 又设 $\mathbb R_n$ 是 $n$ 维实列向量空间，$\mathbb R_n$ 的内积取为标准内积. 令 $\varphi$ 为 $V\to\mathbb R_n$ 的线性映射，它将 $V$ 中向量 $x$ 变为其坐标向量，即若
\[
  x=x_1v_1+x_2v_2+\cdots+x_nv_n,
\]
则 $\varphi(x)=(x_1,x_2,\cdots,x_n)'$. 我们已经知道 $\varphi$ 是同构. 此外若令 $y=y_1v_1+y_2v_2+\cdots+y_nv_n$，则
\[
  (\varphi(x),\varphi(y))=x_1y_1+x_2y_2+\cdots+x_ny_n=(x,y).
\]
因此映射 $\varphi$ 是一个保持内积的同构. 因为向量的长度、向量之间的距离都是由内积决定的，故 $\varphi$ 保持向量的长度和向量之间的距离. 当 $V$ 是酉空间时，类似的结论也成立. 于是我们可以把对抽象内积空间的研究归结为对 $\mathbb R_n$ 或 $\mathbb C_n$ 的研究.

\noindent\textbf{定义 9.4.1}\quad 设 $V$ 与 $U$ 是域 $\mathbb K$ 上的内积空间，$\mathbb K$ 是实数域或复数域，$\varphi$ 是 $V\to U$ 的线性映射. 若对任意的 $x,y\in V$，有
\[
  (\varphi(x),\varphi(y))=(x,y),
\]
则称 $\varphi$ 是 $V\to U$ 的保持内积的线性映射. 又若 $\varphi$ 作为线性映射是同构，则称 $\varphi$ 是内积空间 $V$ 到 $U$ 上的保积同构.

\noindent\textbf{注}\quad (1) 在不引起误解的情况下，我们常把内积空间的保积同构就称为同构.

(2) 保持内积的线性映射一定是单映射，这是因为 $\|\varphi(x)\|=\|x\|$，从 $\|\varphi(x)\|=0$ 得到 $\|x\|=0$，故 $x=0$.

(3) 容易证明保持内积的同构关系是一个等价关系，读者可自己证明之.

我们上面已经提到，若线性映射 $\varphi$ 保持内积，则 $\varphi$ 必保持向量的长度. 现在的问题是如果已知线性映射 $\varphi$ 保持向量的长度或向量之间的距离，那么它是否一定是保积映射？

\noindent\textbf{命题 9.4.1}\quad 若 $\varphi$ 是内积空间 $V$ 到内积空间 $U$ 的保持范数的线性映射，则 $\varphi$ 保持内积.

\noindent\textbf{证明}\quad 向量的范数可以用内积表示，反过来内积也可以用范数来表示. 设 $x,y$ 是 $V$ 中的任意两个向量，则
\[
  \|x+y\|^2=(x+y,x+y)=(x,x)+(y,y)+(x,y)+(y,x),
\]
\[
  \|x-y\|^2=(x-y,x-y)=(x,x)+(y,y)-(x,y)-(y,x),
\]
故
\begin{equation}
  \|x+y\|^2-\|x-y\|^2=2(x,y)+2\overline{(x,y)}.
  \tag{9.4.1}
\end{equation}
另一方面，
\[
  \|x+iy\|^2=(x+iy,x+iy)=(x,x)+(y,y)+i(y,x)-i(x,y),
\]
\[
  \|x-iy\|^2=(x-iy,x-iy)=(x,x)+(y,y)-i(y,x)+i(x,y),
\]
故
\begin{equation}
  \|x+iy\|^2-\|x-iy\|^2=-2i(x,y)+2i\overline{(x,y)}.
  \tag{9.4.2}
\end{equation}
由 (9.4.1) 式和 (9.4.2) 式得
\begin{equation}
  (x,y)=\frac14\|x+y\|^2-\frac14\|x-y\|^2
  +\frac i4\|x+iy\|^2-\frac i4\|x-iy\|^2,
  \tag{9.4.3}
\end{equation}
由此即可得到结论. $\square$

\noindent\textbf{注}\quad 我们仅对复空间进行了讨论，事实上对实空间，可得下列等式：
\begin{equation}
  (x,y)=\frac14\|x+y\|^2-\frac14\|x-y\|^2,
  \tag{9.4.4}
\end{equation}
因此对实空间结论也成立. 由于保持内积与保持范数的等价性，保持内积的同构也称为保范同构或保距同构.

\noindent\textbf{定理 9.4.1}\quad 设 $V$ 与 $U$ 都是 $n$ 维内积空间（同为实空间或同为复空间），若 $\varphi$ 是 $V\to U$ 的线性映射，则下列命题等价：

(1) $\varphi$ 保持内积；

(2) $\varphi$ 是保积同构；

(3) $\varphi$ 将 $V$ 的任一组标准正交基变成 $U$ 的一组标准正交基；

(4) $\varphi$ 将 $V$ 的某一组标准正交基变成 $U$ 的一组标准正交基.

\noindent\textbf{证明}\quad (1) $\Rightarrow$ (2)：$\varphi$ 保持内积，因此 $\varphi$ 为单映射. 由线性映射的维数公式可得
\[
  \dim\operatorname{Im}\varphi=\dim V=\dim U=n,
\]
因此 $\operatorname{Im}\varphi=U$，即 $\varphi$ 是映上的，故为同构.

(2) $\Rightarrow$ (3)：设 $\{e_1,e_2,\cdots,e_n\}$ 是 $V$ 的任意一组标准正交基. 由于 $\varphi$ 保持内积，故对 $i\ne j$，有
\[
  (\varphi(e_i),\varphi(e_j))=(e_i,e_j)=0,
\]
又
\[
  (\varphi(e_i),\varphi(e_i))=(e_i,e_i)=1.
\]
这表明 $\{\varphi(e_1),\varphi(e_2),\cdots,\varphi(e_n)\}$ 是 $U$ 的标准正交基.

(3) $\Rightarrow$ (4)：显然.

(4) $\Rightarrow$ (1)：设 $\{e_1,e_2,\cdots,e_n\}$ 是 $V$ 的标准正交基且 $\{\varphi(e_1),\varphi(e_2),\cdots,\varphi(e_n)\}$ 是 $U$ 的标准正交基. 假设
\[
  u=\sum_{i=1}^n a_ie_i,\quad v=\sum_{i=1}^n b_ie_i,
\]
则
\[
  \varphi(u)=\sum_{i=1}^n a_i\varphi(e_i),\quad
  \varphi(v)=\sum_{i=1}^n b_i\varphi(e_i),
\]
\[
\begin{aligned}
  (\varphi(u),\varphi(v))
  &=
  \left(\sum_{i=1}^n a_i\varphi(e_i),\sum_{i=1}^n b_i\varphi(e_i)\right)\\
  &=a_1\overline{b_1}+a_2\overline{b_2}+\cdots+a_n\overline{b_n}\\
  &=(u,v). \quad \square
\end{aligned}
\]

\noindent\textbf{推论 9.4.1}\quad 两个有限维内积空间 $V$ 与 $U$（同为实空间或同为复空间）同构的充分必要条件是它们有相同的维数.

\noindent\textbf{证明}\quad 只需证明充分性. 设 $\{e_1,e_2,\cdots,e_n\}$ 是 $V$ 的标准正交基，$\{f_1,f_2,\cdots,f_n\}$ 是 $U$ 的标准正交基，令 $\varphi$ 是 $V\to U$ 的线性映射：
\[
  \varphi(e_i)=f_i,\quad i=1,2,\cdots,n,
\]
则 $\varphi$ 将标准正交基变为标准正交基，由定理 9.4.1 可知 $\varphi$ 是保积同构. $\square$

现在我们来讨论同一个内积空间上的保积自同构及其表示矩阵.

\noindent\textbf{定义 9.4.2}\quad 设 $\varphi$ 是内积空间 $V$ 上保持内积的线性变换，若 $V$ 是欧氏空间，则称 $\varphi$ 为正交变换或正交算子；若 $V$ 是酉空间，则称 $\varphi$ 为酉变换或酉算子.

显然正交变换及酉变换都是可逆线性变换. 由定理 9.4.1，正交变换可定义为把欧氏空间中一组标准正交基变成标准正交基的线性变换；酉变换可定义为把酉空间中一组标准正交基变成标准正交基的线性变换.

\noindent\textbf{定理 9.4.2}\quad 设 $\varphi$ 是欧氏空间或酉空间上的线性变换，则 $\varphi$ 是正交变换或酉变换的充分必要条件是 $\varphi$ 非异，且
\[
  \varphi^*=\varphi^{-1}.
\]

\noindent\textbf{证明}\quad 设 $\varphi$ 是欧氏空间 $V$ 上的正交变换，则对 $V$ 中的任意向量 $\alpha,\beta$，有
\[
  (\varphi(\alpha),\beta)=(\varphi(\alpha),\varphi(\varphi^{-1}(\beta)))=(\alpha,\varphi^{-1}(\beta)),
\]
此即 $\varphi^*=\varphi^{-1}$.

反过来，若 $\varphi^*=\varphi^{-1}$，则
\[
  (\varphi(\alpha),\varphi(\beta))=(\alpha,\varphi^*\varphi(\beta))=(\alpha,\beta),
\]
即 $\varphi$ 保持内积，故 $\varphi$ 是正交变换.

对酉变换可类似证明. $\square$

如果在内积空间 $V$ 中取定一组标准正交基，那么正交（酉）变换的表示矩阵有什么特点呢？

\noindent\textbf{定义 9.4.3}\quad 设 $A$ 是 $n$ 阶实方阵，若 $A'=A^{-1}$，则称 $A$ 是正交矩阵. 设 $C$ 是 $n$ 阶复方阵，若 $\overline C'=C^{-1}$，则称 $C$ 是酉矩阵.

\noindent\textbf{定理 9.4.3}\quad 设 $\varphi$ 是欧氏空间（酉空间）$V$ 上的正交变换（酉变换），则在 $V$ 的任一组标准正交基下，$\varphi$ 的表示矩阵是正交矩阵（酉矩阵）.

\noindent\textbf{证明}\quad 由定理 9.4.2，当 $\varphi$ 是正交变换时，若 $\varphi$ 在 $V$ 的一组标准正交基下的表示矩阵为 $A$，则 $\varphi^*$ 在同一组基下的表示矩阵为 $A'$，由 $\varphi^*=\varphi^{-1}$ 得 $A'=A^{-1}$，即 $A$ 是正交矩阵. 同理，当 $\varphi$ 是酉变换时，$\varphi$ 在一组标准正交基下的表示矩阵 $A$ 应适合 $\overline A'=A^{-1}$，即 $A$ 是酉矩阵. $\square$

\noindent\textbf{注}\quad 上述定理的逆命题也是正确的，即若线性变换 $\varphi$ 在一组标准正交基下的表示矩阵为正交（酉）矩阵，则 $\varphi$ 是正交（酉）变换. 这由线性变换与其表示矩阵的关系即得.

由正交矩阵与酉矩阵的定义可知正交矩阵适合条件 $AA'=A'A=I_n$，酉矩阵适合条件 $A\overline A'=\overline A'A=I_n$.

\noindent\textbf{定理 9.4.4}\quad 设 $A=(a_{ij})$ 是 $n$ 阶实矩阵，则 $A$ 是正交矩阵的充分必要条件是：
\[
\begin{gathered}
  a_{i1}a_{j1}+a_{i2}a_{j2}+\cdots+a_{in}a_{jn}=0,\quad i\ne j,\\
  a_{i1}^2+a_{i2}^2+\cdots+a_{in}^2=1,
\end{gathered}
\]
或
\[
\begin{gathered}
  a_{1i}a_{1j}+a_{2i}a_{2j}+\cdots+a_{ni}a_{nj}=0,\quad i\ne j,\\
  a_{1i}^2+a_{2i}^2+\cdots+a_{ni}^2=1.
\end{gathered}
\]
也就是说，$A$ 为正交矩阵的充分必要条件是它的 $n$ 个行向量是 $n$ 维实行向量空间（取标准内积）的标准正交基，或它的 $n$ 个列向量是 $n$ 维实列向量空间（取标准内积）的标准正交基.

\noindent\textbf{证明}\quad 由 $AA'=I_n$ 得到第一个结论，由 $A'A=I_n$ 得到第二个结论. $\square$

同理我们可证明下述定理.

\noindent\textbf{定理 9.4.5}\quad 设 $A=(a_{ij})$ 是 $n$ 阶复矩阵，则 $A$ 是酉矩阵的充分必要条件是：
\[
\begin{gathered}
  a_{i1}\overline{a_{j1}}+a_{i2}\overline{a_{j2}}+\cdots+a_{in}\overline{a_{jn}}=0,\quad i\ne j,\\
  |a_{i1}|^2+|a_{i2}|^2+\cdots+|a_{in}|^2=1,
\end{gathered}
\]
或
\[
\begin{gathered}
  a_{1i}\overline{a_{1j}}+a_{2i}\overline{a_{2j}}+\cdots+a_{ni}\overline{a_{nj}}=0,\quad i\ne j,\\
  |a_{1i}|^2+|a_{2i}|^2+\cdots+|a_{ni}|^2=1.
\end{gathered}
\]
也就是说，$A$ 为酉矩阵的充分必要条件是它的 $n$ 个行向量是 $n$ 维复行向量空间（取标准内积）的标准正交基，或它的 $n$ 个列向量是 $n$ 维复列向量空间（取标准内积）的标准正交基.

\noindent\textbf{定理 9.4.6}\quad 若 $n$ 阶实矩阵 $A$ 是正交矩阵，则

(1) $A$ 的行列式值等于 1 或 $-1$;

(2) $A$ 的特征值的模长等于 1.

\noindent\textbf{证明}\quad (1) 由 $AA'=I_n$，取行列式即得结论.

(2) 设 $\lambda$ 是 $A$ 的特征值，$x$ 是属于 $\lambda$ 的特征向量，则 $Ax=\lambda x$，于是 $\overline x'A'=\overline\lambda\overline x'$. 因此
\[
  \overline x'A'Ax=\overline\lambda\lambda(\overline x'x),
\]
即
\[
  \overline x'x=\overline\lambda\lambda(\overline x'x),
\]
从而 $\overline\lambda\lambda=1$，即 $|\lambda|=1$. $\square$

\noindent\textbf{定理 9.4.7}\quad 若 $n$ 阶复矩阵 $A$ 是酉矩阵，则

(1) $A$ 的行列式值的模长等于 1;

(2) $A$ 的特征值的模长等于 1.

\noindent\textbf{证明}\quad 同上定理. $\square$

\noindent\textbf{例 9.4.1}\quad (1) 单位阵是正交矩阵也是酉矩阵；

(2) 对角阵是正交矩阵的充分必要条件是主对角线上的元素为 1 或 $-1$.

\noindent\textbf{例 9.4.2}\quad (1) 下列矩阵是正交矩阵：
\[
  \left(\begin{array}{cc}
    \cos\theta & -\sin\theta\\
    \sin\theta & \cos\theta
  \end{array}\right);
\]

(2) 下列矩阵为酉矩阵：
\[
  \frac19
  \left(\begin{array}{ccc}
    4+3i & 4i & -6-2i\\
    -4i & 4-3i & -2-6i\\
    6+2i & -2-6i & 1
  \end{array}\right).
\]

下面我们证明矩阵分解理论中的一个重要定理，称为矩阵的 $QR$ 分解.

\noindent\textbf{定理 9.4.8}\quad 设 $A$ 是 $n$ 阶实（复）矩阵，则 $A$ 可分解为
\[
  A=QR,
\]
其中 $Q$ 是正交（酉）矩阵，$R$ 是一个主对角线上的元素均大于等于零的上三角阵，并且若 $A$ 是非异阵，则这样的分解必唯一.

\noindent\textbf{证明}\quad 设 $A$ 是 $n$ 阶实矩阵，$A=(u_1,u_2,\cdots,u_n)$ 是 $A$ 的列分块. 考虑 $n$ 维实列向量空间 $\mathbb R_n$，并取其标准内积，我们先通过类似于 Gram-Schmidt 方法的正交化过程，把 $\{u_1,u_2,\cdots,u_n\}$ 变成一组两两正交的向量 $\{w_1,w_2,\cdots,w_n\}$，并且 $w_k$ 或者是零向量或者是单位向量.

我们用数学归纳法来定义上述向量. 假设 $w_1,\cdots,w_{k-1}$ 已经定义好，现来定义 $w_k$. 令
\[
  v_k=u_k-\sum_{j=1}^{k-1}(u_k,w_j)w_j.
\]
若 $v_k=0$，则令 $w_k=0$；若 $v_k\ne0$，则令 $w_k=\dfrac{v_k}{\|v_k\|}$. 容易验证 $\{w_1,w_2,\cdots,w_n\}$ 是一组两两正交的向量，$w_k$ 或者是零向量或者是单位向量，并且满足
\begin{equation}
  u_k=\sum_{j=1}^{k-1}(u_k,w_j)w_j+\|v_k\|w_k,\quad k=1,2,\cdots,n.
  \tag{9.4.5}
\end{equation}
由 (9.4.5) 式可得
\begin{equation}
  A=(u_1,u_2,\cdots,u_n)=(w_1,w_2,\cdots,w_n)R,
  \tag{9.4.6}
\end{equation}
其中 $R$ 是一个上三角阵且主对角线上的元素依次为 $\|v_1\|,\|v_2\|,\cdots,\|v_n\|$，均大于等于零，并且由 (9.4.5) 式知，如果 $w_k=0$，则 $R$ 的第 $k$ 行元素全为零.

假设 $w_{i_1},w_{i_2},\cdots,w_{i_r}$ 是其中的非零向量全体，由定理 9.2.2 可将它们扩张为 $\mathbb R_n$ 的一组标准正交基 $\{\widetilde w_1,\widetilde w_2,\cdots,\widetilde w_n\}$，其中 $\widetilde w_j=w_j,j=i_1,i_2,\cdots,i_r$. 令 $Q=(\widetilde w_1,\widetilde w_2,\cdots,\widetilde w_n)$，由定理 9.4.4 知 $Q$ 是正交矩阵. 注意到若 $w_k=0$，则 $R$ 的第 $k$ 行元素全为零，此时用 $\widetilde w_k$ 代替 $w_k$ 仍然可使 (9.4.6) 式成立，因此
\[
  A=(u_1,u_2,\cdots,u_n)=(\widetilde w_1,\widetilde w_2,\cdots,\widetilde w_n)R=QR,
\]
从而得到了 $A$ 的 $QR$ 分解.

复矩阵情形的证明完全类似. 至于非异阵 $QR$ 分解的唯一性，我们作为习题留给读者自己证明. $\square$

\noindent\textbf{例 9.4.3}\quad 求下列矩阵的 $QR$ 分解：
\[
  A=
  \left(\begin{array}{ccc}
    1 & 2 & 5\\
    1 & 0 & 1\\
    0 & 1 & 2
  \end{array}\right).
\]

\noindent\textbf{解}\quad 采用与定理 9.4.8 证明中相同的记号，经过计算可得：
\[
\begin{gathered}
  v_1=u_1=(1,1,0)',\quad w_1=\frac1{\sqrt2}(1,1,0)';\\
  v_2=u_2-\sqrt2w_1=(1,-1,1)',\quad w_2=\frac1{\sqrt3}(1,-1,1)';\\
  v_3=u_3-3\sqrt2w_1-2\sqrt3w_2=(0,0,0)',\quad w_3=(0,0,0)'.
\end{gathered}
\]
从而有
\[
  A=(u_1,u_2,u_3)=(w_1,w_2,w_3)
  \left(\begin{array}{ccc}
    \sqrt2 & \sqrt2 & 3\sqrt2\\
    0 & \sqrt3 & 2\sqrt3\\
    0 & 0 & 0
  \end{array}\right).
\]
用 $\widetilde w_3=\dfrac1{\sqrt6}(-1,1,2)'$ 代替 $w_3$ 可得 $A$ 的 $QR$ 分解为
\[
  Q=
  \left(\begin{array}{ccc}
    \dfrac1{\sqrt2} & \dfrac1{\sqrt3} & -\dfrac1{\sqrt6}\\
    \dfrac1{\sqrt2} & -\dfrac1{\sqrt3} & \dfrac1{\sqrt6}\\
    0 & \dfrac1{\sqrt3} & \dfrac2{\sqrt6}
  \end{array}\right),
  \quad
  R=
  \left(\begin{array}{ccc}
    \sqrt2 & \sqrt2 & 3\sqrt2\\
    0 & \sqrt3 & 2\sqrt3\\
    0 & 0 & 0
  \end{array}\right).
\]

\section*{习题 9.4}

1. 证明：正交（酉）变换的积仍是正交（酉）变换，正交（酉）变换的逆也是正交（酉）变换. 正交（酉）矩阵的积仍是正交（酉）矩阵，正交（酉）矩阵的逆仍是正交（酉）矩阵.

2. 证明二阶正交矩阵具有下列形状（按行列式值等于 1 和 $-1$ 进行分类）：
\[
  \left(\begin{array}{cc}
    \cos\theta & -\sin\theta\\
    \sin\theta & \cos\theta
  \end{array}\right),
  \quad
  \left(\begin{array}{cc}
    \cos\theta & \sin\theta\\
    \sin\theta & -\cos\theta
  \end{array}\right).
\]

3. 证明：上三角（或下三角）正交矩阵必是对角阵且主对角线上的元素为 1 或 $-1$. 利用此结论证明定理 9.4.8 中非异阵 $A$ 的 $QR$ 分解必唯一.

4. 求下列矩阵的 $QR$ 分解：
\[
  (1)\quad
  \left(\begin{array}{ccc}
    3 & -1 & 2\\
    0 & 0 & 9\\
    4 & 7 & 11
  \end{array}\right);
  \qquad
  (2)\quad
  \left(\begin{array}{cccc}
    1 & 1 & 0 & 2\\
    0 & 1 & 1 & 1\\
    1 & 1 & 2 & 4\\
    0 & 1 & 1 & 3
  \end{array}\right);
\]
\[
  (3)\quad
  \left(\begin{array}{cccc}
    1 & 2 & 4 & 7\\
    1 & 0 & -2 & -1\\
    1 & 0 & 2 & 3\\
    1 & 2 & 0 & 3
  \end{array}\right).
\]

5. 设 $A,B$ 是 $n$ 阶正交矩阵且 $|A|+|B|=0$，求证：$|A+B|=0$.

6. 求证：正定实对称阵 $A$ 是正交矩阵的充分必要条件是 $A$ 为单位阵.

7. 设 $V,U$ 都是 $n$ 维欧氏空间，$\{e_1,e_2,\cdots,e_n\}$ 和 $\{f_1,f_2,\cdots,f_n\}$ 分别是 $V$ 和 $U$ 的一组基（不一定是标准正交基），线性映射 $\varphi:V\to U$ 满足 $\varphi(e_i)=f_i(i=1,2,\cdots,n)$. 求证：$\varphi$ 是保积同构的充分必要条件是这两组基的 Gram 矩阵相等，即
\[
  G(e_1,e_2,\cdots,e_n)=G(f_1,f_2,\cdots,f_n).
\]

8. 设 $x_1,x_2,\cdots,x_k$ 及 $y_1,y_2,\cdots,y_k$ 是欧氏空间 $V$ 中两组向量. 证明：存在 $V$ 上的正交变换 $\varphi$，使
\[
  \varphi(x_i)=y_i\quad(i=1,2,\cdots,k)
\]
成立的充分必要条件是这两组向量的 Gram 矩阵相等.

9. (1) 设 $e$ 是欧氏空间 $V$ 中长度为 1 的向量，定义线性变换
\[
  \varphi(x)=x-2(e,x)e,
\]
证明：$\varphi$ 是正交变换且 $\det(\varphi)=-1$（上述 $\varphi$ 称为镜像变换）；

(2) 设 $\xi$ 是 $n$ 维欧氏空间 $V$ 中的正交变换，1 是 $\xi$ 的特征值且几何重数等于 $n-1$，证明：必存在 $V$ 中长度为 1 的向量 $e$，使
\[
  \xi(x)=x-2(e,x)e.
\]

10. 设 $n$ 阶矩阵 $M=I_n-2vv'$，其中 $v$ 是 $n$ 维实列向量且 $v'v=1$，这样的 $M$ 称为镜像阵. 设 $\varphi$ 是 $n$ 维欧氏空间 $V$ 上的线性变换，求证：$\varphi$ 是镜像变换的充分必要条件是 $\varphi$ 在 $V$ 的某一组（任一组）标准正交基下的表示矩阵为镜像阵.

11. 设 $u,v$ 是欧氏空间 $V$ 中两个长度相等的不同向量，求证：必存在镜像变换 $\varphi$，使 $\varphi(u)=v$.

12. 证明：$n$ 维欧氏空间 $V$ 中任一正交变换均可表示为不超过 $n$ 个镜像变换之积.

\noindent\textbf{注}\quad 上述结论称为 Cartan-Dieudonn\'e（嘉当--迪厄多内）定理.

\section*{§9.5\quad 自伴随算子}

在下面的几节里我们将讨论与相似标准型类似的问题：一个内积空间上的线性变换适合什么条件，它在一组标准正交基下的表示矩阵具有比较简单的形状？

设 $n$ 维内积空间 $V$ 上的线性变换 $\varphi$ 在一组标准正交基下的表示矩阵为 $A$，在另一组标准正交基下的表示矩阵为 $B$，两组基之间的过渡矩阵为 $P$，我们首先要问：$P$ 是一个什么样的矩阵？

\noindent\textbf{引理 9.5.1}\quad 欧氏空间中两组标准正交基之间的过渡矩阵是正交矩阵，酉空间中两组标准正交基之间的过渡矩阵是酉矩阵.

\noindent\textbf{证明}\quad 设 $V$ 是 $n$ 维欧氏空间，$\{e_1,e_2,\cdots,e_n\}$ 和 $\{f_1,f_2,\cdots,f_n\}$ 是 $V$ 的两组标准正交基且
\[
  \begin{cases}
    f_1=a_{11}e_1+a_{21}e_2+\cdots+a_{n1}e_n,\\
    f_2=a_{12}e_1+a_{22}e_2+\cdots+a_{n2}e_n,\\
    \hspace{3.5em}\cdots\cdots\cdots\cdots\\
    f_n=a_{1n}e_1+a_{2n}e_2+\cdots+a_{nn}e_n.
  \end{cases}
\]
因为 $(f_i,f_i)=1$，故
\begin{equation}
  a_{1i}^2+a_{2i}^2+\cdots+a_{ni}^2=1.
  \tag{9.5.1}
\end{equation}
又若 $i\ne j$，则 $(f_i,f_j)=0$，故
\begin{equation}
  a_{1i}a_{1j}+a_{2i}a_{2j}+\cdots+a_{ni}a_{nj}=0.
  \tag{9.5.2}
\end{equation}
(9.5.1) 式、(9.5.2) 式和定理 9.4.4 表明过渡矩阵
\[
  P=
  \left(\begin{array}{cccc}
    a_{11} & a_{12} & \cdots & a_{1n}\\
    a_{21} & a_{22} & \cdots & a_{2n}\\
    \vdots & \vdots & & \vdots\\
    a_{n1} & a_{n2} & \cdots & a_{nn}
  \end{array}\right)
\]
是正交矩阵. 同理可证明酉空间中两组标准正交基之间的过渡矩阵是酉矩阵. $\square$

根据引理我们知道，当 $V$ 是欧氏空间时，有
\[
  B=P^{-1}AP=P'AP;
\]
当 $V$ 是酉空间时，有
\[
  B=P^{-1}AP=\overline P'AP.
\]

\noindent\textbf{定义 9.5.1}\quad 设 $A,B$ 是 $n$ 阶实矩阵，若存在正交矩阵 $P$，使 $B=P'AP$，则称 $B$ 和 $A$ 正交相似. 设 $A,B$ 是 $n$ 阶复矩阵，若存在酉矩阵 $P$，使 $B=\overline P'AP$，则称 $B$ 和 $A$ 酉相似.

和矩阵的相似关系一样，我们不难证明正交（酉）相似关系是等价关系，即：

(1) $n$ 阶矩阵 $A$ 和自己正交（酉）相似；

(2) 若 $B$ 和 $A$ 正交（酉）相似，则 $A$ 和 $B$ 也正交（酉）相似；

(3) 若 $B$ 和 $A$ 正交（酉）相似，$C$ 和 $B$ 正交（酉）相似，则 $C$ 和 $A$ 也正交（酉）相似.

我们的问题是寻找一类矩阵的正交（酉）相似标准型. 显而易见，正交（酉）相似比普通的相似要求更高，因此对一般的矩阵寻求它们的正交（酉）相似标准型将是很困难的. 在这一节里我们将把注意力限制在一类特殊的矩阵——实对称阵和 Hermite 矩阵上.

\noindent\textbf{定义 9.5.2}\quad 设 $\varphi$ 是内积空间 $V$ 上的线性变换，$\varphi^*$ 是 $\varphi$ 的伴随，若 $\varphi^*=\varphi$，则称 $\varphi$ 是自伴随算子. 当 $V$ 是欧氏空间时，$\varphi$ 也称为对称算子或对称变换；当 $V$ 是酉空间时，$\varphi$ 也称为 Hermite 算子或 Hermite 变换.

\noindent\textbf{例 9.5.1}\quad 设 $V$ 是 $n$ 维内积空间，$V_0$ 是 $V$ 的子空间，$V=V_0\oplus V_0^\perp$. 令 $E$ 是 $V$ 到 $V_0$ 上的正交投影，由命题 9.2.1 知道 $E$ 是自伴随算子.

设 $V$ 是 $n$ 维欧氏空间，$\{e_1,e_2,\cdots,e_n\}$ 是一组标准正交基，若 $\varphi$ 在这组基下的表示矩阵为 $A$，则 $\varphi^*$ 在同一组基下的表示矩阵为 $A'$. 若 $\varphi^*=\varphi$，则
\[
  A'=A.
\]
也就是说欧氏空间上的自伴随算子在任一组标准正交基下的表示矩阵都是实对称阵. 同理可证明酉空间上的自伴随算子在任一组标准正交基下的表示矩阵都是 Hermite 矩阵. 反之亦容易看出，若欧氏空间上的线性变换 $\varphi$ 在某一组标准正交基下的表示矩阵是实对称阵，则 $\varphi^*=\varphi$. 对酉空间也有类似结论.

\noindent\textbf{定理 9.5.1}\quad 设 $V$ 是 $n$ 维酉空间，$\varphi$ 是 $V$ 上的自伴随算子，则 $\varphi$ 的特征值全是实数且属于不同特征值的特征向量互相正交.

\noindent\textbf{证明}\quad 设 $\lambda$ 是 $\varphi$ 的特征值，$x$ 是属于 $\lambda$ 的特征向量，则
\[
\begin{aligned}
  \lambda(x,x)
  &=(\lambda x,x)=(\varphi(x),x)=(x,\varphi^*(x))\\
  &=(x,\varphi(x))=(x,\lambda x)=\overline\lambda(x,x).
\end{aligned}
\]
因为 $(x,x)\ne0$，故 $\overline\lambda=\lambda$，即 $\lambda$ 是实数. 又若设 $\mu$ 是 $\varphi$ 的另一个特征值，$y$ 是属于 $\mu$ 的特征向量，注意到 $\lambda,\mu$ 都是实数，故有
\[
\begin{aligned}
  \lambda(x,y)
  &=(\lambda x,y)=(\varphi(x),y)=(x,\varphi^*(y))\\
  &=(x,\varphi(y))=(x,\mu y)=\mu(x,y).
\end{aligned}
\]
由于 $\lambda\ne\mu$，故 $(x,y)=0$，即 $x\perp y$. $\square$

\noindent\textbf{推论 9.5.1}\quad Hermite 矩阵的特征值全是实数，实对称阵的特征值也全是实数. 这两种矩阵属于不同特征值的特征向量互相正交.

\noindent\textbf{证明}\quad Hermite 矩阵的结论是定理的显然推论，而实对称阵也是 Hermite 矩阵，因此结论成立. $\square$

\noindent\textbf{定理 9.5.2}\quad 设 $V$ 是 $n$ 维内积空间，$\varphi$ 是 $V$ 上的自伴随算子，则存在 $V$ 的一组标准正交基，使 $\varphi$ 在这组基下的表示矩阵为实对角阵，且这组基恰为 $\varphi$ 的 $n$ 个线性无关的特征向量.

\noindent\textbf{证明}\quad 首先需要说明的是，若 $V$ 是欧氏空间，则由于自伴随算子 $\varphi$ 的特征值都是实数，故有实的特征向量. 不妨设 $u$ 是 $\varphi$ 的特征向量，令 $v_1=\dfrac{u}{\|u\|}$，则 $v_1$ 是 $\varphi$ 的长度等于 1 的特征向量. 我们对维数 $n$ 用归纳法.

若 $\dim V=1$，结论已成立. 设对小于 $n$ 维的内积空间结论成立. 令 $W$ 为由 $v_1$ 张成的子空间，$W^\perp$ 为 $W$ 的正交补空间，则 $W$ 是 $\varphi$ 的不变子空间且
\[
  V=W\oplus W^\perp,\quad \dim W^\perp=n-1.
\]
由命题 9.3.1 可知 $W^\perp$ 是 $\varphi^*=\varphi$ 的不变子空间. 将 $\varphi$ 限制在 $W^\perp$ 上仍是自伴随算子. 由归纳假设，存在 $W^\perp$ 的一组标准正交基 $\{v_2,\cdots,v_n\}$，使 $\varphi$ 在这组基下的表示矩阵为实对角阵，且 $\{v_2,\cdots,v_n\}$ 是其特征向量. 因此，$\{v_1,v_2,\cdots,v_n\}$ 构成了 $V$ 的一组标准正交基，$\varphi$ 在这组基下的表示矩阵为实对角阵，且 $\{v_1,v_2,\cdots,v_n\}$ 为 $\varphi$ 的 $n$ 个线性无关的特征向量. $\square$

\noindent\textbf{定理 9.5.3}\quad (1) 设 $A$ 是 $n$ 阶实对称阵，则存在正交矩阵 $P$，使 $P'AP$ 为对角阵，且 $P$ 的 $n$ 个列向量恰为 $A$ 的 $n$ 个两两正交的单位特征向量.

(2) 设 $A$ 是 $n$ 阶 Hermite 矩阵，则存在酉矩阵 $P$，使 $\overline P'AP$ 为实对角阵，且 $P$ 的 $n$ 个列向量恰为 $A$ 的 $n$ 个两两正交的单位特征向量.

\noindent\textbf{证明}\quad 由定理 9.5.2 即知实对称阵正交相似于对角阵，Hermite 矩阵酉相似于实对角阵. 从 §6.2 知道 $P$ 的列向量都是 $A$ 的特征向量，又 $P$ 是正交（酉）矩阵，故这些列向量两两正交且长度等于 1. $\square$

上述对角阵称为实对称（Hermite）矩阵 $A$ 的正交（酉）相似标准型. 对角阵主对角线上的元素就是 $A$ 的全体特征值.

现在的问题是实对称（Hermite）矩阵在正交（酉）相似关系下的全系不变量是什么呢？

\noindent\textbf{推论 9.5.2}\quad 实对称阵的全体特征值是实对称阵在正交相似关系下的全系不变量，Hermite 矩阵的全体特征值是 Hermite 矩阵在酉相似关系下的全系不变量.

\noindent\textbf{证明}\quad 只证明实的情形. 显然正交相似的矩阵有相同的特征值. 另一方面，由定理 9.5.3 知道只需对对角阵证明，若它们的特征值相同，则必正交相似即可. 设
\[
  B=\operatorname{diag}\{\lambda_1,\lambda_2,\cdots,\lambda_n\},\quad
  D=\operatorname{diag}\{\lambda_{i_1},\lambda_{i_2},\cdots,\lambda_{i_n}\},
\]
其中 $\{\lambda_{i_1},\lambda_{i_2},\cdots,\lambda_{i_n}\}$ 是 $\{\lambda_1,\lambda_2,\cdots,\lambda_n\}$ 的一个排列. 由于任一排列可通过若干次对换来实现，因此只要证明对 $B$ 的第 $(i,i)$ 元素和第 $(j,j)$ 元素对换后得到的矩阵与 $B$ 正交相似即可. 设 $P_{ij}$ 是第一类初等矩阵，则 $P_{ij}$ 是正交矩阵且 $P_{ij}'=P_{ij}$，因此 $P_{ij}'BP_{ij}$ 和 $B$ 正交相似. 这就是我们要证明的结论. $\square$

实对称阵的正交相似标准型理论在实二次型理论中有重要的应用. 从几何的角度来看，以三维欧氏空间为例，就是要在不改变量度的条件下将二次曲面的方程化为标准型. 这也是解析几何中所谓的主轴问题.

\noindent\textbf{定理 9.5.4}\quad 设 $f(x)=x'Ax$ 是 $n$ 元实二次型，系数矩阵 $A$ 的特征值为 $\lambda_1,\lambda_2,\cdots,\lambda_n$，则 $f$ 经过正交变换 $x=Py$ 可以化为下列标准型：
\[
  \lambda_1y_1^2+\lambda_2y_2^2+\cdots+\lambda_ny_n^2.
\]
因此，$f$ 的正惯性指数等于 $A$ 的正特征值的个数，负惯性指数等于 $A$ 的负特征值的个数，$f$ 的秩等于 $A$ 的非零特征值的个数.

\noindent\textbf{证明}\quad 注意到正交相似既是相似又是合同，故由定理 9.5.3 即得结论. $\square$

\noindent\textbf{推论 9.5.3}\quad 设 $f(x)=x'Ax$ 是 $n$ 元实二次型，则 $f$ 是正定型当且仅当系数矩阵 $A$ 的特征值全是正数，$f$ 是负定型当且仅当 $A$ 的特征值全是负数，$f$ 是半正定型当且仅当 $A$ 的特征值全非负，$f$ 是半负定型当且仅当 $A$ 的特征值全非正.

下面我们通过具体例子来说明对实对称阵 $A$，如何求正交矩阵 $P$，使 $P'AP$ 是对角阵. 我们的方法和 §6.2 类似，只是增加了特征向量的标准正交化过程.

\noindent\textbf{例 9.5.2}\quad 求正交矩阵 $P$，使 $P'AP$ 为对角阵，其中
\[
  A=
  \left(\begin{array}{ccc}
    4 & 2 & 2\\
    2 & 4 & 2\\
    2 & 2 & 4
  \end{array}\right).
\]

\noindent\textbf{解}\quad 先求特征值
\[
  |\lambda I-A|=
  \left|\begin{array}{ccc}
    \lambda-4 & -2 & -2\\
    -2 & \lambda-4 & -2\\
    -2 & -2 & \lambda-4
  \end{array}\right|
  =(\lambda-8)(\lambda-2)^2.
\]
因此，$A$ 的特征值为 $\lambda_1=8,\lambda_2=\lambda_3=2$. 当 $\lambda=8$ 时，求解齐次线性方程组 $(\lambda I-A)x=0$，得到基础解系（只有一个向量）：
\[
  \eta_1=(1,1,1)'.
\]
当 $\lambda=2$ 时，求解齐次线性方程组 $(\lambda I-A)x=0$，得到基础解系（有两个向量）：
\[
  \eta_2=(-1,1,0)',\quad \eta_3=(-1,0,1)'.
\]
由于实对称阵属于不同特征值的特征向量必正交，因此只要对上面两个向量正交化，用 Gram-Schmidt 方法将 $\eta_2,\eta_3$ 正交化得到
\[
  \xi_2=(-1,1,0)',\quad \xi_3=\left(-\frac12,-\frac12,1\right)'.
\]
再将 $\eta_1,\xi_2,\xi_3$ 化为单位向量得到
\[
  v_1=\left(\frac1{\sqrt3},\frac1{\sqrt3},\frac1{\sqrt3}\right)',\quad
  v_2=\left(-\frac1{\sqrt2},\frac1{\sqrt2},0\right)',\quad
  v_3=\left(-\frac1{\sqrt6},-\frac1{\sqrt6},\frac2{\sqrt6}\right)'.
\]
令
\[
  P=(v_1,v_2,v_3)=
  \left(\begin{array}{ccc}
    \dfrac1{\sqrt3} & -\dfrac1{\sqrt2} & -\dfrac1{\sqrt6}\\
    \dfrac1{\sqrt3} & \dfrac1{\sqrt2} & -\dfrac1{\sqrt6}\\
    \dfrac1{\sqrt3} & 0 & \dfrac2{\sqrt6}
  \end{array}\right),
\]
于是
\[
  P'AP=
  \left(\begin{array}{ccc}
    8 & 0 & 0\\
    0 & 2 & 0\\
    0 & 0 & 2
  \end{array}\right).
\]

\noindent\textbf{注}\quad 上述方法也适用于 Hermite 矩阵，即对 Hermite 矩阵 $A$，求酉矩阵 $P$，使 $\overline P'AP$ 为对角阵.

\noindent\textbf{例 9.5.3}\quad 设 $A$ 是三阶实对称阵，$A$ 的特征值为 $0,3,3$. 已知属于特征值 $0$ 的特征向量为 $v_1=(1,1,1)'$，又向量 $v_2=(-1,1,0)'$ 是属于特征值 $3$ 的特征向量，求矩阵 $A$.

\noindent\textbf{解}\quad 设 $A$ 属于特征值 $3$ 的另一特征向量 $v_3=(x_1,x_2,x_3)'$ 和 $v_1,v_2$ 都正交，则
\[
  \begin{cases}
    x_1+x_2+x_3=0,\\
    -x_1+x_2=0,
  \end{cases}
\]
求出一个非零解 $v_3=(1,1,-2)'$. 因为 $v_1,v_2,v_3$ 已经两两正交，故只需将 $v_1,v_2,v_3$ 标准化，得到
\[
  \xi_1=
  \left(\begin{array}{c}
    \dfrac1{\sqrt3}\\
    \dfrac1{\sqrt3}\\
    \dfrac1{\sqrt3}
  \end{array}\right),\quad
  \xi_2=
  \left(\begin{array}{c}
    -\dfrac1{\sqrt2}\\
    \dfrac1{\sqrt2}\\
    0
  \end{array}\right),\quad
  \xi_3=
  \left(\begin{array}{c}
    \dfrac1{\sqrt6}\\
    \dfrac1{\sqrt6}\\
    -\dfrac2{\sqrt6}
  \end{array}\right).
\]
令
\[
  P=(\xi_1,\xi_2,\xi_3)=
  \left(\begin{array}{ccc}
    \dfrac1{\sqrt3} & -\dfrac1{\sqrt2} & \dfrac1{\sqrt6}\\
    \dfrac1{\sqrt3} & \dfrac1{\sqrt2} & \dfrac1{\sqrt6}\\
    \dfrac1{\sqrt3} & 0 & -\dfrac2{\sqrt6}
  \end{array}\right),\quad
  B=
  \left(\begin{array}{ccc}
    0 & 0 & 0\\
    0 & 3 & 0\\
    0 & 0 & 3
  \end{array}\right),
\]
则
\[
  A=PBP'=
  \left(\begin{array}{ccc}
    2 & -1 & -1\\
    -1 & 2 & -1\\
    -1 & -1 & 2
  \end{array}\right).
\]

\noindent\textbf{例 9.5.4}\quad 设 $A$ 是 $n$ 阶实对称阵且 $A^3=I_n$，证明：$A=I_n$.

\noindent\textbf{证明}\quad 设 $A$ 的特征值为 $\lambda$，则 $\lambda^3=1$. 因为 $\lambda$ 是实数，故 $\lambda=1$，这就是说 $A$ 的特征值全是 $1$. 由定理 9.5.3，存在正交矩阵 $P$，使 $P'AP=I_n$，于是 $A=PI_nP'=I_n$. $\square$

\section*{习题 9.5}

1. 求正交矩阵 $P$，使 $P'AP$ 为对角阵：
\[
\begin{array}{ll}
(1)\ A=
\left(\begin{array}{ccc}
2 & 2 & -2\\
2 & 5 & -4\\
-2 & -4 & 5
\end{array}\right);
&
(2)\ A=
\left(\begin{array}{ccc}
5 & -1 & 3\\
-1 & 5 & -3\\
3 & -3 & 3
\end{array}\right);
\\[3em]
(3)\ A=
\left(\begin{array}{rrrr}
-1 & -3 & 3 & -3\\
-3 & -1 & -3 & 3\\
3 & -3 & -1 & -3\\
-3 & 3 & -3 & -1
\end{array}\right);
&
(4)\ A=
\left(\begin{array}{rrrr}
3 & 1 & 0 & -1\\
1 & 3 & 0 & 1\\
0 & 0 & 4 & 0\\
-1 & 1 & 0 & 3
\end{array}\right).
\end{array}
\]

2. 求酉矩阵 $P$，使 $\overline P'AP$ 为对角阵：
\[
\begin{array}{ll}
(1)\ A=
\left(\begin{array}{cc}
3 & 2+2i\\
2-2i & 1
\end{array}\right);
&
(2)\ A=
\left(\begin{array}{cc}
3 & 2-i\\
2+i & 7
\end{array}\right).
\end{array}
\]

3. 设实对称阵
\[
  A=
  \left(\begin{array}{rrrr}
    a & 1 & 1 & -1\\
    1 & a & -1 & 1\\
    1 & -1 & a & 1\\
    -1 & 1 & 1 & a
  \end{array}\right)
\]
有一个单特征值 $-3$，求 $a$ 的值并求正交矩阵 $P$，使 $P'AP$ 为对角阵.

4. 设三阶实对称阵 $A$ 的各行元素之和均为 $3$，向量 $\alpha_1=(-1,2,-1)'$，$\alpha_2=(0,-1,1)'$ 是线性方程组 $Ax=0$ 的两个解.

(1) 求 $A$ 的特征值和特征向量；

(2) 求正交矩阵 $Q$ 和对角阵 $D$，使 $Q'AQ=D$.

5. 设四阶实对称阵 $A$ 的特征值为 $0,0,0,4$，且属于特征值 $0$ 的线性无关特征向量为 $(-1,1,0,0)'$，$(-1,0,1,0)'$，$(-1,0,0,1)'$，求出矩阵 $A$.

6. 设实二次型 $f(x_1,x_2,x_3)=2x_1^2+5x_2^2+5x_3^2+2ax_1x_2+2bx_1x_3-8x_2x_3$，经过正交变换后 $f$ 可化为 $y_1^2+y_2^2+cy_3^2$，求出 $a,b,c$ 的值和正交变换矩阵 $P$.

7. 已知曲面 $x^2+ay^2+z^2+2bxy+2xz+2yz=4$ 可以经过正交变换
\[
  \left(\begin{array}{c}
    x\\ y\\ z
  \end{array}\right)
  =P
  \left(\begin{array}{c}
    u\\ v\\ w
  \end{array}\right)
\]
化为椭圆柱面 $v^2+4w^2=4$，求 $a,b$ 之值和正交矩阵 $P$.

8. 设 $\varphi,\psi$ 是内积空间 $V$ 上的两个自伴随算子，求证：

(1) $\varphi+\psi,\varphi\psi+\psi\varphi,i(\varphi\psi-\psi\varphi)$ 也都是自伴随算子；

(2) $\varphi\psi$ 是自伴随算子的充分必要条件是 $\varphi\psi=\psi\varphi$.

9. 设 $V$ 是 $n$ 维欧氏空间，$G$ 是 $V$ 关于基 $\{e_1,e_2,\cdots,e_n\}$ 的度量矩阵，线性变换 $\varphi$ 在这组基下的表示矩阵为 $A$，求 $\varphi$ 是自伴随算子的充分必要条件.

10. 设 $\varphi,\psi$ 是 $n$ 维内积空间 $V$ 上的两个自伴随算子，证明：$V$ 有一组由 $\varphi$ 和 $\psi$ 的公共特征向量构成的标准正交基的充分必要条件是 $\varphi\psi=\psi\varphi$.

\section*{§ 9.6\quad 复正规算子}

我们现在来讨论这样一个问题：如果 $n$ 维酉空间 $V$ 上的线性变换 $\varphi$ 在一组标准正交基下的表示矩阵是对角阵，则 $\varphi$ 必须满足什么条件？设 $\{e_1,e_2,\cdots,e_n\}$ 是 $V$ 的标准正交基，$\varphi$ 在这组基下的表示矩阵为
\[
  A=\operatorname{diag}\{c_1,c_2,\cdots,c_n\},
\]
则
\[
  \varphi(e_i)=c_ie_i,\quad i=1,2,\cdots,n,
\]
即 $\varphi$ 以 $c_1,c_2,\cdots,c_n$ 为特征值，$e_1,e_2,\cdots,e_n$ 为对应的特征向量. 设 $\varphi^*$ 为 $\varphi$ 的伴随，则 $\varphi^*$ 在这组基下的表示矩阵为
\[
  \overline A'=\operatorname{diag}\{\overline c_1,\overline c_2,\cdots,\overline c_n\}.
\]
于是，我们有
\[
  \overline A'A=A\overline A',\quad \varphi\varphi^*=\varphi^*\varphi.
\]

\noindent\textbf{定义 9.6.1}\quad 设 $\varphi$ 是内积空间 $V$ 上的线性变换，$\varphi^*$ 是其伴随，若 $\varphi\varphi^*=\varphi^*\varphi$，则称 $\varphi$ 是 $V$ 上的正规算子. 为了不引起混淆，我们也称酉空间（欧氏空间）$V$ 上的正规算子 $\varphi$ 为复正规算子（实正规算子）. 复矩阵 $A$ 若适合 $\overline A'A=A\overline A'$，则称其为复正规矩阵. 实矩阵 $A$ 若适合 $A'A=AA'$，则称其为实正规矩阵.

\noindent\textbf{注}\quad (1) 容易验证酉算子（酉矩阵）和 Hermite 算子（Hermite 矩阵）都是复正规算子（矩阵）. 正交变换（正交矩阵）和对称变换（实对称矩阵）都是实正规算子（矩阵）. 我们将在下一节中详细讨论实正规算子和实正规矩阵.

(2) 容易证明酉空间（欧氏空间）$V$ 上的线性变换 $\varphi$ 是复（实）正规算子的充分必要条件是 $\varphi$ 在 $V$ 的某一组或任一组标准正交基下的表示矩阵都是复（实）正规矩阵. 因此，复（实）矩阵的正规性在酉（正交）相似下是不变的.

从上面的论证我们看出：一个复矩阵如果酉相似于对角阵，则它一定是复正规矩阵，那么复正规矩阵是否必酉相似于对角阵呢？

我们先证明复正规算子（复正规矩阵）的特征值和特征向量的一些重要性质.

\noindent\textbf{引理 9.6.1}\quad 设 $\varphi$ 是内积空间 $V$ 上的正规算子，则对任意的 $\alpha\in V$，成立
\[
  \|\varphi(\alpha)\|=\|\varphi^*(\alpha)\|.
\]

\noindent\textbf{证明}\quad 由 $\varphi$ 的正规性，有
\[
\begin{aligned}
  \|\varphi(\alpha)\|^2
  &= (\varphi(\alpha),\varphi(\alpha))=(\alpha,\varphi^*\varphi(\alpha))\\
  &= (\alpha,\varphi\varphi^*(\alpha))=(\varphi^*(\alpha),\varphi^*(\alpha))\\
  &= \|\varphi^*(\alpha)\|^2.
\end{aligned}
\]

\noindent\textbf{命题 9.6.1}\quad 设 $V$ 是 $n$ 维酉空间，$\varphi$ 是 $V$ 上的正规算子.

(1) 向量 $u$ 是 $\varphi$ 属于特征值 $\lambda$ 的特征向量的充分必要条件为 $u$ 是 $\varphi^*$ 属于特征值 $\overline\lambda$ 的特征向量；

(2) 属于 $\varphi$ 不同特征值的特征向量必正交.

\noindent\textbf{证明}\quad (1) 若 $\lambda$ 是任一数，则 $(\lambda I-\varphi)^*=\overline\lambda I-\varphi^*$，且
\[
  (\lambda I-\varphi)(\overline\lambda I-\varphi^*)=(\overline\lambda I-\varphi^*)(\lambda I-\varphi),
\]
即 $\lambda I-\varphi$ 也是正规算子. 于是由引理 9.6.1，
\[
  \|(\lambda I-\varphi)(\alpha)\|=\|(\overline\lambda I-\varphi^*)(\alpha)\|
\]
对一切 $\alpha\in V$ 成立，故 $(\lambda I-\varphi)(u)=0$ 当且仅当 $(\overline\lambda I-\varphi^*)(u)=0$ 成立.

(2) 设 $\varphi(u)=\lambda u,\varphi(v)=\mu v$ 且 $\lambda\ne\mu$，则由 (1) 知 $\varphi^*(v)=\overline\mu v$，于是
\[
  \lambda(u,v)=(\lambda u,v)=(\varphi(u),v)=(u,\varphi^*(v))=(u,\overline\mu v)=\mu(u,v).
\]
因为 $\lambda\ne\mu$，故 $(u,v)=0$. $\square$

\noindent\textbf{引理 9.6.2}\quad 设 $V$ 是 $n$ 维酉空间，$\varphi$ 是 $V$ 上的线性变换，又 $\{e_1,e_2,\cdots,e_n\}$ 是 $V$ 的一组标准正交基. 设 $\varphi$ 在这组基下的表示矩阵 $A$ 是一个上三角阵，则 $\varphi$ 是正规算子的充分必要条件是 $A$ 为对角阵.

\noindent\textbf{证明}\quad 若 $A$ 是对角阵，则 $A\overline A'=\overline A'A$，故 $\varphi\varphi^*=\varphi^*\varphi$，即 $\varphi$ 是正规算子. 反之，设 $\varphi$ 是正规算子. 由于 $A$ 是上三角阵，可记 $A=(a_{ij})$，$a_{ij}=0\ (i>j)$. 于是 $\varphi(e_1)=a_{11}e_1$，再由上面的命题可知 $\varphi^*(e_1)=\overline a_{11}e_1$. 另一方面，有
\[
  \varphi^*(e_1)=\overline a_{11}e_1+\overline a_{12}e_2+\cdots+\overline a_{1n}e_n.
\]
因此 $a_{1j}=0$ 对一切 $j>1$ 成立. 又因为 $A$ 是上三角，所以
\[
  \varphi(e_2)=a_{22}e_2,
\]
故又有 $\varphi^*(e_2)=\overline a_{22}e_2$ 及 $a_{2j}=0\ (j>2)$. 不断这样做下去即得 $A$ 是对角阵. $\square$

\noindent\textbf{定理 9.6.1 (Schur（舒尔）定理)}\quad 设 $V$ 是 $n$ 维酉空间，$\varphi$ 是 $V$ 上的线性算子，则存在 $V$ 的一组标准正交基，使 $\varphi$ 在这组基下的表示矩阵为上三角阵.

\noindent\textbf{证明}\quad 对 $V$ 的维数 $n$ 用数学归纳法. 当 $n=1$ 时结论显然成立. 设对 $n-1$ 维酉空间结论成立，现证 $n$ 维酉空间的情形. 由于 $V$ 是复线性空间，故 $\varphi^*$ 总存在特征值与特征向量，即有
\[
  \varphi^*(e)=\lambda e.
\]
设 $W$ 是由 $e$ 张成的一维子空间的正交补空间，由命题 9.3.1 知 $W$ 是 $(\varphi^*)^*=\varphi$ 的不变子空间，将 $\varphi$ 限制在 $W$ 上得到 $W$ 上的一个线性变换. 注意到 $\dim W=n-1$，故由归纳假设，存在 $W$ 的一组标准正交基 $\{e_1,e_2,\cdots,e_{n-1}\}$，使 $\varphi|_W$ 在这组基下的表示矩阵为上三角阵. 令 $e_n=\dfrac{e}{\|e\|}$，则 $\{e_1,e_2,\cdots,e_n\}$ 成为 $V$ 的一组标准正交基，使 $\varphi$ 在这组基下的表示矩阵为上三角阵. $\square$

\noindent\textbf{推论 9.6.1 (Schur 定理)}\quad 任一 $n$ 阶复矩阵均酉相似于一个上三角阵.

利用上面两个命题，我们立即得到下列定理.

\noindent\textbf{定理 9.6.2}\quad 设 $V$ 是 $n$ 维酉空间，$\varphi$ 是 $V$ 上的线性算子，则 $\varphi$ 为正规算子的充分必要条件是存在 $V$ 的一组标准正交基，使 $\varphi$ 在这组基下的表示矩阵是对角阵，特别，这组基恰为 $\varphi$ 的 $n$ 个线性无关的特征向量.

\noindent\textbf{定理 9.6.3}\quad 复矩阵 $A$ 为复正规矩阵的充分必要条件是 $A$ 酉相似于对角阵.

显而易见，复正规矩阵的特征值就是复正规矩阵在酉相似关系下的全系不变量，即两个复正规矩阵酉相似的充分必要条件是它们具有相同的特征值.

由推论 6.2.3 我们知道，对一个可对角化线性变换，它的所有特征子空间的直和就是全空间 $V$. 对复正规算子我们也有类似的结论.

\noindent\textbf{命题 9.6.2}\quad 设 $\varphi$ 是 $n$ 维酉空间 $V$ 上的线性算子，$\lambda_1,\lambda_2,\cdots,\lambda_k$ 是 $\varphi$ 的全体不同特征值，$V_1,V_2,\cdots,V_k$ 是对应的特征子空间，则 $\varphi$ 是正规算子的充分必要条件是
\begin{equation}
  V=V_1\perp V_2\perp\cdots\perp V_k.
  \tag{9.6.1}
\end{equation}

\noindent\textbf{证明}\quad 设 $\varphi$ 是正规算子，则它是一个可对角化线性变换，因此
\[
  V=V_1\oplus V_2\oplus\cdots\oplus V_k.
\]
又从命题 9.6.1 知道，若 $i\ne j$，则 $V_i\perp V_j$，所以 (9.6.1) 式成立.

反之，若 (9.6.1) 式成立，则在每个 $V_i$ 中取一组标准正交基，将这些基向量组成 $V$ 的一组标准正交基. 因为每个 $V_i$ 都是 $\varphi$ 的特征子空间，即 $\varphi(\alpha)=\lambda_i\alpha$ 对一切 $\alpha\in V_i$ 成立，故 $\varphi$ 在这组基下的表示矩阵是对角阵，因此 $\varphi$ 是正规算子. $\square$

由于酉矩阵是正规矩阵，故有下列结论.

\noindent\textbf{定理 9.6.4}\quad 任一 $n$ 阶酉矩阵必酉相似于下列对角阵：
\[
  \operatorname{diag}\{c_1,c_2,\cdots,c_n\},
\]
其中 $c_i$ 为模长等于 $1$ 的复数.

\noindent\textbf{证明}\quad 由定理 9.6.3 知酉矩阵酉相似于 $\operatorname{diag}\{c_1,c_2,\cdots,c_n\}$. 由于与酉矩阵酉相似的矩阵仍是酉矩阵，故 $\operatorname{diag}\{c_1,c_2,\cdots,c_n\}$ 是酉矩阵，因此 $|c_i|=1$. $\square$

\section*{习题 9.6}

1. 构造一个二阶复矩阵，它可对角化但不是复正规矩阵.

2. 设复矩阵 $A$ 是斜 Hermite 矩阵，即 $\overline A'=-A$. 证明：$A$ 必酉相似于对角阵
\[
  \operatorname{diag}\{c_1,c_2,\cdots,c_n\},
\]
其中 $c_i$ 是零或纯虚数.

3. 求酉矩阵 $P$，使 $\overline P'AP$ 是对角阵，其中
\[
  A=
  \left(\begin{array}{cc}
    1 & i\\
    i & 1
  \end{array}\right).
\]

4. 设 $A$ 是复正规矩阵，求证：$A$ 是 Hermite 矩阵的充分必要条件是 $A$ 的特征值全是实数.

5. 设 $A$ 是复正规矩阵，求证：$A$ 是酉矩阵的充分必要条件是 $A$ 的特征值全是模长等于 $1$ 的复数.

6. 设 $A$ 是复正规矩阵，求证：$A$ 是幂零阵（即存在正整数 $k$，使 $A^k=O$）的充分必要条件是 $A=O$.

7. 设 $\varphi$ 是 $n$ 维酉空间 $V$ 上的正规算子，求证：对 $\varphi$ 的任一特征值 $\lambda_0$，总有
\[
  V=\operatorname{Ker}(\varphi-\lambda_0 I_V)\perp \operatorname{Im}(\varphi-\lambda_0 I_V).
\]

8. 设 $\lambda_1,\lambda_2,\cdots,\lambda_n$ 是 $n$ 阶复矩阵 $A=(a_{ij})$ 的特征值，求证：
\[
  \sum_{i=1}^n|\lambda_i|^2\le \sum_{i,j=1}^n|a_{ij}|^2,
\]
等号成立的充分必要条件是 $A$ 为复正规矩阵.

\section*{* § 9.7\quad 实正规矩阵}

由上一节我们知道，适合 $AA'=A'A$ 的实矩阵 $A$ 称为正规矩阵. 实正规矩阵的正交相似标准型比复正规矩阵的酉相似标准型要复杂一些，这是因为任一复矩阵总有复特征值与复特征向量，而实正规矩阵可能没有实特征值及实特征向量. 为了求得实正规矩阵的正交相似标准型，我们采用“几何”的方法.

首先利用欧氏空间 $V$ 上的正规算子 $\varphi$ 的极小多项式的不可约分解可将 $V$ 分解为若干个不变子空间的正交直和，并且 $\varphi$ 限制在每个不变子空间上的极小多项式不超过二次. 这样，问题就归结为研究极小多项式次数不超过二次的正规算子. 因为极小多项式为一次的线性变换就是数量变换，所以下去就对极小多项式为二次不可约多项式的正规算子进行讨论. 这样就可以得到实正规矩阵的正交相似标准型.

\noindent\textbf{引理 9.7.1}\quad 设 $V$ 是 $n$ 维欧氏空间，$f(x)$ 是一个实系数多项式，若 $\varphi$ 是 $V$ 上的正规算子，则 $f(\varphi)$ 也是 $V$ 上的正规算子.

\noindent\textbf{证明}\quad 设
\[
  f(x)=a_0+a_1x+\cdots+a_mx^m,
\]
则
\[
  f(\varphi)=a_0I+a_1\varphi+\cdots+a_m\varphi^m,\quad
  f(\varphi)^*=a_0I+a_1\varphi^*+\cdots+a_m(\varphi^*)^m=f(\varphi^*).
\]
由 $\varphi\varphi^*=\varphi^*\varphi$，不难验证
\[
  f(\varphi)f(\varphi)^*=f(\varphi)^*f(\varphi).
\]

\noindent\textbf{引理 9.7.2}\quad 设 $\varphi$ 是欧氏空间 $V$ 上的正规算子，$f(x),g(x)$ 是互素的实系数多项式. 假设 $u\in\operatorname{Ker} f(\varphi),v\in\operatorname{Ker} g(\varphi)$，则
\[
  (u,v)=0.
\]

\noindent\textbf{证明}\quad 因为 $f(x),g(x)$ 互素，故存在实系数多项式 $s(x),t(x)$，使
\[
  f(x)s(x)+g(x)t(x)=1,
\]
于是
\[
  f(\varphi)s(\varphi)+g(\varphi)t(\varphi)=I.
\]
因此 $u=g(\varphi)t(\varphi)(u)$，
\[
  (u,v)=(g(\varphi)t(\varphi)(u),v)=(t(\varphi)(u),g(\varphi)^*(v)).
\]
由上面的引理知 $g(\varphi)$ 是正规算子且 $g(\varphi)(v)=0$，再由引理 9.6.1 得 $g(\varphi)^*(v)=0$，因此 $(u,v)=0$. $\square$

\noindent\textbf{定理 9.7.1}\quad 设 $V$ 是 $n$ 维欧氏空间，$\varphi$ 是 $V$ 上的正规算子. 令 $g(x)$ 是 $\varphi$ 的极小多项式，$g_1(x),\cdots,g_k(x)$ 为 $g(x)$ 的所有互不相同的首一不可约因式，且
\[
  W_i=\operatorname{Ker} g_i(\varphi)\quad (i=1,\cdots,k),
\]
则

(1) $g(x)=g_1(x)\cdots g_k(x)$，其中 $\deg g_i(x)\le 2$；

(2) $V=W_1\perp\cdots\perp W_k$；

(3) $W_i$ 是 $\varphi$ 的不变子空间，用 $\varphi_i$ 表示 $\varphi$ 在 $W_i$ 上的限制，则 $\varphi_i$ 是 $W_i$ 上的正规算子且 $g_i(x)$ 是 $\varphi_i$ 的极小多项式.

\noindent\textbf{证明}\quad (1) 设 $\varphi$ 在 $V$ 的一组标准正交基下的表示矩阵为正规矩阵 $A$，则 $A$ 也是复正规矩阵，因此 $A$ 酉相似于对角阵，于是 $A$ 的极小多项式 $g(x)$ 在复数域上无重根，从而在实数域上无重因式. 又不可约实系数多项式的次数小于等于 $2$，故 (1) 的结论成立.

(2) 令 $f_i(x)=g(x)/g_i(x)$，则 $f_1(x),\cdots,f_k(x)$ 互素. 由 §5.3 习题 7 知，存在实系数多项式 $h_1(x),\cdots,h_k(x)$，使
\[
  f_1(x)h_1(x)+\cdots+f_k(x)h_k(x)=1.
\]
对任意的 $v\in V$，
\begin{equation}
  v=f_1(\varphi)h_1(\varphi)(v)+\cdots+f_k(\varphi)h_k(\varphi)(v).
  \tag{9.7.1}
\end{equation}
注意到 $g_i(\varphi)f_i(\varphi)=g(\varphi)=0$，故对任一 $i$，$f_i(\varphi)h_i(\varphi)(v)\in\operatorname{Ker} g_i(\varphi)=W_i$，于是由 (9.7.1) 式知
\[
  V=W_1+\cdots+W_k.
\]
当 $i\ne j$ 时，$g_i(x)$ 与 $g_j(x)$ 互素，由引理 9.7.2 得到 $W_i\perp W_j$，故
\[
  V=W_1\perp\cdots\perp W_k.
\]

(3) 任取 $u\in W_i=\operatorname{Ker}g_i(\varphi)$，注意到 $\varphi\varphi^*=\varphi^*\varphi$，故
\[
  g_i(\varphi)(\varphi(u))=\varphi g_i(\varphi)(u)=0,\quad
  g_i(\varphi)(\varphi^*(u))=\varphi^*g_i(\varphi)(u)=0,
\]
于是 $\varphi(u)\in W_i$ 且 $\varphi^*(u)\in W_i$，因此 $W_i$ 既是 $\varphi$ 的不变子空间，也是 $\varphi^*$ 的不变子空间. 由伴随的定义容易验证 $\varphi_i=\varphi|_{W_i}$ 的伴随 $\varphi_i^*$ 等于 $\varphi^*|_{W_i}$，于是
\[
  \varphi_i\varphi_i^*=\varphi|_{W_i}\varphi^*|_{W_i}
  =\varphi^*|_{W_i}\varphi|_{W_i}=\varphi_i^*\varphi_i,
\]
即 $\varphi_i$ 是 $W_i$ 上的正规算子. 由于 $W_i=\operatorname{Ker}g_i(\varphi)$，故 $0=g_i(\varphi)|_{W_i}=g_i(\varphi|_{W_i})=g_i(\varphi_i)$，即 $\varphi_i$ 适合多项式 $g_i(x)$，于是 $\varphi_i$ 的极小多项式 $m_i(x)\mid g_i(x)$. 注意到 $g_i(x)$ 首一不可约，故只能是 $m_i(x)=g_i(x)$. $\square$

接下去我们要讨论极小多项式是二次不可约多项式的实正规算子的表示矩阵，先对比较简单的情形即极小多项式为 $x^2+1$ 的正规算子进行讨论，再过渡到一般情形.

\noindent\textbf{引理 9.7.3}\quad 设 $V$ 是 $n$ 维欧氏空间，$\varphi$ 是 $V$ 上的正规算子且 $\varphi$ 适合多项式 $g(x)=x^2+1$. 设 $v\in V,u=\varphi(v)$，则
\[
  \varphi^*(v)=-u,\quad \varphi^*(u)=v,
\]
且 $\|u\|=\|v\|,u\perp v$.

\noindent\textbf{证明}\quad 由 $\varphi(u)=\varphi^2(v)=-v$，得
\[
\begin{aligned}
  0
  &= \|\varphi(v)-u\|^2+\|\varphi(u)+v\|^2\\
  &= \|\varphi(v)\|^2-2(\varphi(v),u)+\|u\|^2
     +\|\varphi(u)\|^2+2(\varphi(u),v)+\|v\|^2.
\end{aligned}
\]
因为 $\varphi$ 是正规算子，由引理 9.6.1 知 $\|\varphi(v)\|^2=\|\varphi^*(v)\|^2$，$\|\varphi(u)\|^2=\|\varphi^*(u)\|^2$. 于是
\[
\begin{aligned}
  0
  &= \|\varphi^*(v)\|^2+2(\varphi^*(v),u)+\|u\|^2
     +\|\varphi^*(u)\|^2-2(\varphi^*(u),v)+\|v\|^2\\
  &= \|\varphi^*(v)+u\|^2+\|\varphi^*(u)-v\|^2.
\end{aligned}
\]
从而 $\varphi^*(v)=-u,\varphi^*(u)=v$. 又
\[
  (v,u)=(\varphi^*(u),u)=(u,\varphi(u))=(u,-v)=-(v,u),
\]
因此 $(v,u)=0$，即 $v\perp u$. 最后
\[
  \|v\|^2=(v,v)=(\varphi^*(u),v)=(u,\varphi(v))=(u,u)=\|u\|^2.
\]

\noindent\textbf{引理 9.7.4}\quad 设 $V$ 是 $n$ 维欧氏空间，$\varphi$ 是 $V$ 上的正规算子且 $\varphi$ 适合多项式 $g(x)=(x-a)^2+b^2$，其中 $a,b$ 都是实数且 $b\ne0$. 设 $v\in V,u=b^{-1}(\varphi-aI)(v)$，则 $\|u\|=\|v\|,u\perp v$，且
\[
  \varphi(v)=av+bu,\quad \varphi(u)=-bv+au;
\]
\[
  \varphi^*(v)=av-bu,\quad \varphi^*(u)=bv+au.
\]

\noindent\textbf{证明}\quad 令 $\psi=b^{-1}(\varphi-aI)$，则 $\psi$ 是 $V$ 上的正规算子且适合多项式 $x^2+1$. 又 $u=\psi(v),\psi^*=b^{-1}(\varphi^*-aI)$，故由引理 9.7.3 即得所需结论. $\square$

\noindent\textbf{定理 9.7.2}\quad 设 $\varphi$ 是 $n$ 维欧氏空间 $V$ 上的正规算子，$\varphi$ 的极小多项式为
\[
  g(x)=(x-a)^2+b^2,
\]
其中 $a,b$ 是实数且 $b\ne0$，则存在 $V$ 的 $s$ 个二维子空间 $V_1,\cdots,V_s$，使
\[
  V=V_1\perp\cdots\perp V_s,
\]
且每个 $V_i$ 都有标准正交基 $\{u_i,v_i\}$，满足
\begin{equation}
  \varphi(u_i)=au_i-bv_i,\quad \varphi(v_i)=bu_i+av_i.
  \tag{9.7.2}
\end{equation}

\noindent\textbf{证明}\quad 对 $V$ 的维数 $n$ 进行归纳. 当 $n=0$ 时，结论是平凡的. 当 $n=1$ 时，任取 $V$ 中的非零向量 $v$，设 $\varphi(v)=cv$，其中 $c$ 是实数，则 $c$ 是 $\varphi$ 的特征值. 因此 $c$ 必须适合 $\varphi$ 的极小多项式 $g(x)$，但 $g(c)=(c-a)^2+b^2>0$，矛盾. 这说明 $n=1$ 的情形不可能发生. 设维数小于 $n$ 时结论已成立，现证 $n$ 维欧氏空间的情形.

任取 $V$ 中长度等于 $1$ 的向量 $v_1$，令 $u_1=b^{-1}(\varphi-aI)(v_1)$，则 $u_1,v_1$ 是两个长度等于 $1$ 的正交向量. 令 $V_1$ 是由 $u_1,v_1$ 张成的子空间，则由引理 9.7.4 知
\[
  \varphi(u_1)=au_1-bv_1,\quad \varphi(v_1)=bu_1+av_1,
\]
\[
  \varphi^*(u_1)=au_1+bv_1,\quad \varphi^*(v_1)=-bu_1+av_1,
\]
因此 $V_1$ 是 $\varphi$ 和 $\varphi^*$ 的不变子空间. 令 $W=V_1^\perp$，则由命题 9.3.1 知 $W$ 是 $\varphi$ 和 $\varphi^*$ 的不变子空间. 考虑 $\varphi$ 在 $W$ 上的限制，容易验证 $\varphi|_W$ 是 $W$ 上的正规算子且极小多项式仍为 $g(x)$. 因为 $\dim W=n-2$，故由归纳假设，存在 $s-1$ 个二维子空间 $V_2,\cdots,V_s$，使
\[
  W=V_2\perp\cdots\perp V_s,
\]
且每个 $V_i\ (i=2,\cdots,s)$ 都有标准正交基 $\{u_i,v_i\}$ 满足 (9.7.2) 式. 因此
\[
  V=V_1\perp V_2\perp\cdots\perp V_s,
\]
且每个 $V_i\ (i=1,2,\cdots,s)$ 都有满足条件的标准正交基 $\{u_i,v_i\}$，结论得证. 特别，$\varphi$ 在 $V$ 的标准正交基 $\{u_1,v_1,u_2,v_2,\cdots,u_s,v_s\}$ 下的表示矩阵为分块对角阵：
\[
  A=\operatorname{diag}\left\{
  \left(\begin{array}{cc}
    a & b\\
    -b & a
  \end{array}\right),
  \left(\begin{array}{cc}
    a & b\\
    -b & a
  \end{array}\right),
  \cdots,
  \left(\begin{array}{cc}
    a & b\\
    -b & a
  \end{array}\right)
  \right\}.
\]

将上面的讨论综合起来，就可得到这一节的主要结论.

\noindent\textbf{定理 9.7.3}\quad 设 $V$ 是 $n$ 维欧氏空间，$\varphi$ 是 $V$ 上的正规算子，则存在一组标准正交基，使 $\varphi$ 在这组基下的表示矩阵为下列分块对角阵：
\begin{equation}
  \operatorname{diag}\{A_1,\cdots,A_r,c_{2r+1},\cdots,c_n\},
  \tag{9.7.3}
\end{equation}
其中 $c_j\ (j=2r+1,\cdots,n)$ 是实数，$A_i$ 为形如
\[
  \left(\begin{array}{cc}
    a_i & b_i\\
    -b_i & a_i
  \end{array}\right)
\]
的二阶实矩阵.

\noindent\textbf{证明}\quad 由定理 9.7.1 知
\[
  V=W_1\perp W_2\perp\cdots\perp W_k,
\]
其中 $W_i=\operatorname{Ker}g_i(\varphi)$，$g_i(x)$ 是次数不超过 $2$ 的多项式，并且 $\varphi$ 在 $W_i$ 上的限制 $\varphi_i$ 是 $W_i$ 上的正规算子，其极小多项式是 $g_i(x)$. 对每个 $W_i$，若 $g_i(x)$ 是二次多项式，则由定理 9.7.2 知 $W_i$ 可分解为若干个二维子空间的正交直和. 若 $g_i(x)=x-c_i$，则 $\varphi_i-c_iI=0$，即 $\varphi_i=c_iI$. 由此即可得到所需结论. $\square$

\noindent\textbf{注}\quad (9.7.3) 式就是实正规矩阵的正交相似标准型. 在不计主对角线上分块的次序的意义下，实正规矩阵的正交相似标准型是唯一确定的. 因此，实正规矩阵的全体特征值是正交相似关系的全系不变量.

利用实正规矩阵的正交相似标准型，我们立即可得到正交矩阵的正交相似标准型.

\noindent\textbf{定理 9.7.4}\quad 设 $A$ 是 $n$ 阶正交矩阵，则 $A$ 正交相似于下列分块对角阵：
\[
  \operatorname{diag}\{A_1,\cdots,A_r;1,\cdots,1;-1,\cdots,-1\},
\]
其中
\[
  A_i=
  \left(\begin{array}{cc}
    \cos\theta_i & \sin\theta_i\\
    -\sin\theta_i & \cos\theta_i
  \end{array}\right),\quad i=1,\cdots,r.
\]

\noindent\textbf{证明}\quad 由于正交矩阵是正规矩阵，故由定理 9.7.3 知道 $A$ 必正交相似于
\[
  \operatorname{diag}\{A_1,\cdots,A_r;c_{2r+1},\cdots,c_n\}.
\]
又由正交矩阵的性质知 $|c_j|=1\ (j=2r+1,\cdots,n)$. 另一方面，设
\[
  A_i=
  \left(\begin{array}{cc}
    a_i & b_i\\
    -b_i & a_i
  \end{array}\right),
\]
则 $a_i^2+b_i^2=1$，故可设 $a_i=\cos\theta_i,b_i=\sin\theta_i$. $\square$

设 $A$ 是 $n$ 阶实方阵，若 $A$ 适合 $A'=-A$，则称 $A$ 是实反对称阵或实斜对称阵. 由于 $AA'=-A^2=A'A$，故实反对称阵是正规矩阵. 若 $P$ 是正交矩阵，则
\[
  (P'AP)'=P'A'P=-P'AP.
\]
因此，反对称性在正交相似下保持不变. 下面是实反对称阵的正交相似标准型.

\noindent\textbf{定理 9.7.5}\quad 设 $A$ 是实反对称阵，则 $A$ 正交相似于下列分块对角阵：
\[
  \operatorname{diag}\{B_1,\cdots,B_r;0,\cdots,0\},
\]
其中
\[
  B_i=
  \left(\begin{array}{cc}
    0 & b_i\\
    -b_i & 0
  \end{array}\right),\quad i=1,\cdots,r.
\]

\noindent\textbf{证明}\quad 设 $A$ 正交相似于
\[
  \operatorname{diag}\{B_1,\cdots,B_r;c_{2r+1},\cdots,c_n\}.
\]
由 $B'=-B$ 即得 $c_j=0\ (j=2r+1,\cdots,n)$. 再由 $B_i'=-B_i$ 知道，$B_i$ 必具有定理所需之形状. $\square$

\noindent\textbf{推论 9.7.1}\quad 实反对称阵的秩必是偶数，且其实特征值必为 $0$，虚特征值为纯虚数.

\section*{习题 9.7}

1. 设 $\varphi$ 是 $n$ 维欧氏空间 $V$ 上的正规算子，其极小多项式为 $g(x)=(x-a)^2+b^2$，其中 $b\ne0$，证明：$\varphi$ 是 $V$ 上的自同构且
\[
  \varphi^*=(a^2+b^2)\varphi^{-1}.
\]

2. 设 $\varphi$ 是 $n$ 维欧氏空间 $V$ 上的正规算子，$\psi$ 是 $V$ 上的线性算子，满足 $\varphi\psi=\psi\varphi$，证明：
\[
  \varphi^*\psi=\psi\varphi^*.
\]

3. 设 $A,B$ 是 $n$ 阶实正规矩阵，求证：若 $A,B$ 相似，则它们必正交相似.

4. 设 $\varphi$ 是 $n$ 维欧氏空间 $V$ 上的正交变换，若 $\det\varphi=1$，则称 $\varphi$ 是一个旋转；若 $\det\varphi=-1$，则称 $\varphi$ 是一个反射. 求证：

(1) 奇数维空间的旋转必有保持不动的非零向量，即存在 $0\ne v\in V$，使 $\varphi(v)=v$；

(2) 反射必有反向的非零向量，即存在 $0\ne v\in V$，使 $\varphi(v)=-v$.

5. 证明：$n$ 阶实方阵必正交相似于下列分块上三角阵：
\[
  \left(\begin{array}{ccccc}
    A_1 & & & & *\\
    & \ddots & & &\\
    & & A_r & &\\
    & & & c_1 &\\
    & & & & \ddots\\
    & & & & c_k
  \end{array}\right),
\]
其中 $A_i\ (i=1,\cdots,r)$ 是二阶实矩阵，其特征值为 $a_i\pm b_i\sqrt{-1}$，$c_j\ (j=1,\cdots,k)$ 是实数.

6. 设 $A,B$ 是实方阵且分块矩阵
\[
  \left(\begin{array}{cc}
    A & C\\
    O & B
  \end{array}\right)
\]
是实正规矩阵，求证：$C=O$ 且 $A,B$ 也是实正规矩阵.

7. 利用第 5 题和第 6 题证明实正规矩阵的正交相似标准型定理.

8. 设 $A,B$ 是 $n$ 阶实正规矩阵，满足 $AB=BA$，求证：存在 $n$ 阶正交矩阵 $P$，使 $P'AP$ 与 $P'BP$ 同时为正交相似标准型.

\section*{* § 9.8\quad 谱分解与极分解}

设 $V$ 是 $n$ 维酉空间，$\varphi$ 是 $V$ 上的正规算子，$\lambda_1,\lambda_2,\cdots,\lambda_k$ 为 $\varphi$ 的全体不同特征值. 由命题 9.6.2 知道，存在 $V$ 的一组标准正交基，使 $\varphi$ 在这组基下的表示矩阵为对角阵
\[
  A=\operatorname{diag}\{\lambda_1,\cdots,\lambda_1;\lambda_2,\cdots,\lambda_2;\cdots;\lambda_k,\cdots,\lambda_k\}.
\]
假设特征值 $\lambda_i$ 的重数等于 $r_i$，则在分解式
\[
  V=W_1\perp W_2\perp\cdots\perp W_k
\]
中，$W_i$ 是 $\varphi$ 的特征子空间且维数等于 $r_i$. 若记 $D_i$ 为如下对角阵
\[
  D_i=\operatorname{diag}\{0,\cdots,0;\cdots;1,\cdots,1;\cdots;0,\cdots,0\},
\]
即主对角元有 $r_i$ 个 $1$，其余都是 $0$ 的对角阵，则
\[
  A=\lambda_1D_1+\lambda_2D_2+\cdots+\lambda_kD_k.
\]
诸 $D_i$ 满足条件 $D_i^2=D_i,D_iD_j=O\ (i\ne j),D_1+D_2+\cdots+D_k=I_n$.

对实对称阵也有类似的结论. 如果把上面的结论“翻译”成几何语言就是下面的谱分解定理. 我们将给出一个用“几何”方法的证明.

\noindent\textbf{定理 9.8.1}\quad 设 $V$ 是有限维内积空间，$\varphi$ 是 $V$ 上的线性算子，当 $V$ 是酉空间时 $\varphi$ 为正规算子；当 $V$ 是欧氏空间时 $\varphi$ 为自伴随算子. 设 $\lambda_1,\lambda_2,\cdots,\lambda_k$ 是 $\varphi$ 的全体不同特征值，$W_i$ 为 $\varphi$ 属于 $\lambda_i$ 的特征子空间，则 $V$ 是 $W_1,W_2,\cdots,W_k$ 的正交直和. 设 $E_i$ 是 $V$ 到 $W_i$ 上的正交投影，则 $\varphi$ 有下列分解式：
\begin{equation}
  \varphi=\lambda_1E_1+\lambda_2E_2+\cdots+\lambda_kE_k.
  \tag{9.8.1}
\end{equation}

\noindent\textbf{证明}\quad 由命题 9.6.2 知道
\[
  V=W_1\perp W_2\perp\cdots\perp W_k.
\]
又因为 $E_i$ 是 $V\to W_i$ 的正交投影，故
\[
  I=E_1+E_2+\cdots+E_k,
\]
注意 $\varphi E_i=\lambda_iE_i$，于是
\[
\begin{aligned}
  \varphi
  &= \varphi E_1+\varphi E_2+\cdots+\varphi E_k\\
  &= \lambda_1E_1+\lambda_2E_2+\cdots+\lambda_kE_k.
\end{aligned}
\]

谱分解定理有许多重要的应用，下面我们择要介绍其应用.

\noindent\textbf{引理 9.8.1}\quad 设
\[
  f_j(x)=\prod_{i\ne j}\frac{x-\lambda_i}{\lambda_j-\lambda_i},
\]
则 $E_j=f_j(\varphi)$.

\noindent\textbf{证明}\quad 当 $i\ne j$ 时，$E_iE_j=0$，故
\[
  \varphi^2=\lambda_1^2E_1+\lambda_2^2E_2+\cdots+\lambda_k^2E_k.
\]
同理不难证明
\[
  \varphi^n=\lambda_1^nE_1+\lambda_2^nE_2+\cdots+\lambda_k^nE_k
\]
对一切正整数 $n$ 成立. 若设
\[
  f(x)=a_0+a_1x+\cdots+a_nx^n,
\]
则
\[
\begin{aligned}
  f(\varphi)
  &= a_0I+a_1\varphi+\cdots+a_n\varphi^n\\
  &= a_0\left(\sum_{i=1}^kE_i\right)
     +a_1\left(\sum_{i=1}^k\lambda_iE_i\right)
     +\cdots
     +a_n\left(\sum_{i=1}^k\lambda_i^nE_i\right)\\
  &= \sum_{i=1}^kf(\lambda_i)E_i.
\end{aligned}
\]
由 $f_j(\lambda_j)=1,f_j(\lambda_i)=0\ (i\ne j)$ 即得 $f_j(\varphi)=E_j$. $\square$

\noindent\textbf{推论 9.8.1}\quad 设 $\varphi$ 是酉空间 $V$ 上的线性算子，则 $\varphi$ 是正规算子的充分必要条件是存在复系数多项式 $f(x)$，使 $\varphi^*=f(\varphi)$.

\noindent\textbf{证明}\quad 若存在复系数多项式 $f(x)$，使 $\varphi^*=f(\varphi)$，则 $\varphi\varphi^*=\varphi^*\varphi$，即 $\varphi$ 是正规算子. 若 $\varphi$ 是正规算子，由定理 9.8.1 知 $\varphi$ 存在谱分解：
\[
  \varphi=\lambda_1E_1+\lambda_2E_2+\cdots+\lambda_kE_k.
\]
注意到 $E_i$ 是自伴随算子，故
\[
  \varphi^*=\overline\lambda_1E_1+\overline\lambda_2E_2+\cdots+\overline\lambda_kE_k.
\]
采用与引理 9.8.1 相同的记号，令 $f(x)=\sum_{j=1}^k\overline\lambda_j f_j(x)$，则
\[
  f(\varphi)=\sum_{j=1}^k\overline\lambda_j f_j(\varphi)=\sum_{j=1}^k\overline\lambda_j E_j=\varphi^*.
\]

\noindent\textbf{定义 9.8.1}\quad 设 $\varphi$ 是内积空间 $V$ 上的自伴随算子，若对任意的非零向量 $\alpha\in V$，总有 $(\varphi(\alpha),\alpha)>0\ ((\varphi(\alpha),\alpha)\ge0)$，则称 $\varphi$ 为正定（半正定）自伴随算子.

容易证明，$\varphi$ 是正定自伴随算子当且仅当 $\varphi$ 在 $V$ 的一组标准正交基下的表示矩阵是正定 Hermite 矩阵（酉空间时）或正定实对称阵（欧氏空间时）；$\varphi$ 是半正定自伴随算子当且仅当 $\varphi$ 在 $V$ 的一组标准正交基下的表示矩阵是半正定 Hermite 矩阵（酉空间时）或半正定实对称阵（欧氏空间时）.

虽然下面的定理用标准型很容易证明，但是用谱分解来证明亦不失为一个好方法.

\noindent\textbf{定理 9.8.2}\quad 设 $\varphi$ 是酉空间 $V$ 上的正规算子. 若 $\varphi$ 的特征值全是实数，则 $\varphi$ 是自伴随算子；若 $\varphi$ 的特征值全是非负实数，则 $\varphi$ 是半正定自伴随算子；若 $\varphi$ 的特征值全是正实数，则 $\varphi$ 是正定自伴随算子；若 $\varphi$ 的特征值的模长等于 $1$，则 $\varphi$ 是酉算子.

\noindent\textbf{证明}\quad 设 $\varphi$ 的谱分解为
\[
  \varphi=\lambda_1E_1+\lambda_2E_2+\cdots+\lambda_kE_k,
\]
则
\[
  \varphi^*=\overline\lambda_1E_1+\overline\lambda_2E_2+\cdots+\overline\lambda_kE_k.
\]
若 $\varphi$ 的特征值全是实数，则 $\varphi^*=\varphi$，即 $\varphi$ 是自伴随算子. 若 $\lambda_i$ 全是非负实数，则对任意的非零向量 $\alpha\in V$，有
\[
  \alpha=E_1(\alpha)+E_2(\alpha)+\cdots+E_k(\alpha),
\]
\[
  \varphi(\alpha)=\lambda_1E_1(\alpha)+\lambda_2E_2(\alpha)+\cdots+\lambda_kE_k(\alpha),
\]
从而
\[
  (\varphi(\alpha),\alpha)=\lambda_1\|E_1(\alpha)\|^2+\lambda_2\|E_2(\alpha)\|^2+\cdots+\lambda_k\|E_k(\alpha)\|^2\ge0.
\]
同理，若特征值全是正实数，则 $\varphi$ 是正定自伴随算子. 最后，若 $|\lambda_i|=1$，则
\[
\begin{aligned}
  \varphi\varphi^*
  &= \lambda_1\overline\lambda_1E_1+\lambda_2\overline\lambda_2E_2+\cdots+\lambda_k\overline\lambda_kE_k\\
  &= |\lambda_1|^2E_1+|\lambda_2|^2E_2+\cdots+|\lambda_k|^2E_k\\
  &= E_1+E_2+\cdots+E_k=I,
\end{aligned}
\]
即 $\varphi$ 是酉算子. $\square$

下面的定理也可用代数方法证明，但是唯一性的证明会比较困难，而用谱分解定理证明则比较容易.

\noindent\textbf{定理 9.8.3}\quad 设 $V$ 是有限维内积空间，$\varphi$ 是 $V$ 上的半正定自伴随算子，则存在 $V$ 上唯一的半正定自伴随算子 $\psi$，使 $\psi^2=\varphi$.

\noindent\textbf{证明}\quad 设 $\varphi$ 的谱分解为
\[
  \varphi=\lambda_1E_1+\lambda_2E_2+\cdots+\lambda_kE_k.
\]
令 $d_i=\sqrt{\lambda_i}\ (i=1,2,\cdots,k)$，则
\[
  \psi=d_1E_1+d_2E_2+\cdots+d_kE_k
\]
适合 $\psi^2=\varphi$ 且 $\psi$ 也是半正定自伴随算子.

现设 $\theta$ 是 $V$ 上的半正定自伴随算子且 $\theta^2=\varphi$，我们要证明 $\theta=\psi$. 令
\[
  \theta=b_1F_1+b_2F_2+\cdots+b_rF_r
\]
是 $\theta$ 的谱分解，其中 $F_i$ 是正交投影算子且 $b_i$ 为非负实数. 由 $\theta^2=\varphi$ 得
\[
  \varphi=b_1^2F_1+b_2^2F_2+\cdots+b_r^2F_r.
\]
因为非负实数 $b_i$ 互不相同，故 $b_1^2,b_2^2,\cdots,b_r^2$ 是 $\varphi$ 的全体不同特征值，于是 $r=k$，$b_i^2=\lambda_i$，从而 $b_i=d_i$（这里允许差一个次序）. 注意到 $E_i(V)$ 及 $F_i(V)$ 都是 $\varphi$ 的关于特征值 $\lambda_i$ 的特征子空间，因此 $F_i=E_i$，这就证明了 $\theta=\psi$. $\square$

\noindent\textbf{推论 9.8.2}\quad 设 $A$ 是半正定实对称（Hermite）矩阵，则必存在唯一的半正定实对称（Hermite）矩阵 $B$，使 $A=B^2$.

\noindent\textbf{定理 9.8.4}\quad 设 $V$ 是 $n$ 维酉空间（欧氏空间），$\varphi$ 是 $V$ 上的任一线性算子，则存在 $V$ 上的酉算子（正交算子）$\omega$ 以及 $V$ 上的半正定自伴随算子 $\psi$，使 $\varphi=\omega\psi$，其中 $\psi$ 是唯一的，并且若 $\varphi$ 是非异线性算子，则 $\omega$ 也唯一.

\noindent\textbf{证明}\quad 若已有 $\varphi=\omega\psi$，其中 $\omega$ 为酉算子（正交算子），$\psi$ 为半正定自伴随算子，则
\[
  \varphi^*=\psi^*\omega^*=\psi\omega^*,
\]
\[
  \varphi^*\varphi=\psi\omega^*\omega\psi=\psi^2.
\]
由定义容易验证 $\varphi^*\varphi$ 是半正定自伴随算子，故由定理 9.8.3 知，$\psi$ 被 $\varphi$ 唯一确定.

现来证明存在性. 令 $\psi$ 是定理 9.8.3 中的 $\psi$ 且使 $\psi^2=\varphi^*\varphi$. 若 $\varphi$ 是非异线性变换，则 $\psi$ 也是非异的. 事实上，这时
\begin{equation}
  (\psi(v),\psi(v))=(\psi^2(v),v)=(\varphi^*\varphi(v),v)=(\varphi(v),\varphi(v))
  \tag{9.8.2}
\end{equation}
对一切 $v\in V$ 成立. 显然从 $\varphi$ 非异即不难推出 $\psi$ 也非异. 这时可令 $\omega=\varphi\psi^{-1}$，只需证明 $\omega$ 是酉算子（正交算子）即可. 注意到
\[
  \omega^*=(\varphi\psi^{-1})^*=(\psi^{-1})^*\varphi^*=(\psi^*)^{-1}\varphi^*=\psi^{-1}\varphi^*,
\]
\[
  \omega^*\omega=\psi^{-1}\varphi^*\varphi\psi^{-1}=\psi^{-1}\psi^2\psi^{-1}=I,
\]
因此 $\omega$ 是酉算子（正交算子）.

若 $\varphi$ 不是非异线性变换，现来定义 $\omega$. 设 $W=\operatorname{Im}\psi$，$W^\perp$ 是其正交补空间. 定义 $W=\operatorname{Im}\psi\to\operatorname{Im}\varphi$ 的映射 $\eta$ 如下：若 $\psi(u)\in W$，则
\[
  \eta(\psi(u))=\varphi(u).
\]
这时我们必须验证 $\eta$ 定义的合理性，即若 $\psi(u)=\psi(v)$，则必须有 $\varphi(u)=\varphi(v)$. 由 (9.8.2) 式可知对任意的 $\alpha\in V$，$\|\psi(\alpha)\|^2=\|\varphi(\alpha)\|^2$，因此 $\psi(u-v)=0$ 当且仅当 $\varphi(u-v)=0$. 这表明 $\eta$ 是一个合理定义的映射. 又 $\eta$ 显然是线性的.

再定义 $W^\perp\to(\operatorname{Im}\varphi)^\perp$ 的映射 $\xi$：由 (9.8.2) 式可知 $\varphi$ 与 $\psi$ 的核空间相同，故像空间的维数相等. 于是 $W^\perp$ 与 $(\operatorname{Im}\varphi)^\perp$ 的维数相等，故必存在一个保持内积的同构，这个同构记为 $\xi$. 由于 $V=W\oplus W^\perp$，故 $V$ 中任一向量 $v$ 均可唯一地表示为
\[
  v=w+w',
\]
其中 $w\in W,w'\in W^\perp$. 令
\[
  \omega(v)=\eta(w)+\xi(w'),
\]
不难看出 $\omega$ 是 $V$ 上的线性变换. 若设 $w=\psi(u)$，则
\[
\begin{aligned}
  (\omega(v),\omega(v))
  &= (\eta(w)+\xi(w'),\eta(w)+\xi(w'))\\
  &= (\varphi(u)+\xi(w'),\varphi(u)+\xi(w'))\\
  &= (\varphi(u),\varphi(u))+(\xi(w'),\xi(w'))\\
  &= (\psi(u),\psi(u))+(w',w')\\
  &= (w,w)+(w',w')\\
  &= (v,v),
\end{aligned}
\]
因此 $\omega$ 是酉算子（正交算子）. 显然这时 $\varphi=\omega\psi$. $\square$

我们称 $\varphi=\omega\psi$ 是一个极分解. $\varphi$ 也可作这样的分解：$\varphi=\psi_1\omega$，这里 $\omega$ 仍为之前的酉算子（正交算子），$\psi_1=\omega\psi\omega^*$ 为半正定自伴随算子. 这个式子也称为极分解.

\noindent\textbf{推论 9.8.3}\quad 设 $A$ 是 $n$ 阶实矩阵，则存在 $n$ 阶正交矩阵 $Q$ 以及 $n$ 阶半正定实对称阵 $S$，使 $A=QS$. 设 $B$ 是 $n$ 阶复矩阵，则存在 $n$ 阶酉矩阵 $U$ 以及 $n$ 阶半正定 Hermite 矩阵 $H$，使 $B=UH$. 当 $A,B$ 为非异阵时，上述分解式被唯一确定.

矩阵极分解的另一形式读者不难自己写出. 从定理 9.8.4 的证明过程很容易得到非异阵极分解的计算方法，我们将通过下面的例子进行阐述. 对于奇异阵的极分解，若将定理 9.8.4 的证明过程转化为计算方法则过于繁琐，因此我们将通过 §9.9 矩阵的奇异值分解来求相应的极分解.

\noindent\textbf{例 9.8.1}\quad 求下列非异阵的极分解：
\[
  A=
  \left(\begin{array}{rrr}
    10 & 7 & 11\\
    5 & 10 & 5\\
    -5 & -1 & 2
  \end{array}\right).
\]

\noindent\textbf{解}\quad 经过计算可得
\[
  A'A=
  \left(\begin{array}{ccc}
    150 & 125 & 125\\
    125 & 150 & 125\\
    125 & 125 & 150
  \end{array}\right).
\]
采用与例 9.5.2 相同的计算方法可得正交矩阵
\[
  P=
  \left(\begin{array}{ccc}
    -\dfrac1{\sqrt2} & -\dfrac1{\sqrt6} & \dfrac1{\sqrt3}\\
    \dfrac1{\sqrt2} & -\dfrac1{\sqrt6} & \dfrac1{\sqrt3}\\
    0 & \dfrac2{\sqrt6} & \dfrac1{\sqrt3}
  \end{array}\right),
\]
使 $P'A'AP=\operatorname{diag}\{25,25,400\}$. 令
\[
  S=P
  \left(\begin{array}{ccc}
    5 & 0 & 0\\
    0 & 5 & 0\\
    0 & 0 & 20
  \end{array}\right)
  P'
  =
  \left(\begin{array}{ccc}
    10 & 5 & 5\\
    5 & 10 & 5\\
    5 & 5 & 10
  \end{array}\right),
\]
则 $S$ 为正定阵且 $A'A=S^2$. 再令
\[
  Q=AS^{-1}=
  \left(\begin{array}{rrr}
    10 & 7 & 11\\
    5 & 10 & 5\\
    -5 & -1 & 2
  \end{array}\right)
  \left(\begin{array}{ccc}
    10 & 5 & 5\\
    5 & 10 & 5\\
    5 & 5 & 10
  \end{array}\right)^{-1}
  =
  \left(\begin{array}{ccc}
    \dfrac35 & 0 & \dfrac45\\
    0 & 1 & 0\\
    -\dfrac45 & 0 & \dfrac35
  \end{array}\right),
\]
则 $A=QS$ 即为所求的极分解.

\section*{习题 9.8}

1. 设 $\varphi$ 是 $n$ 维酉空间 $V$ 上的线性算子，证明：$\varphi$ 正规的充分必要条件是 $\varphi=\varphi_1+i\varphi_2$，其中 $\varphi_1,\varphi_2$ 为自伴随算子且 $\varphi_1\varphi_2=\varphi_2\varphi_1$.

2. 设 $\varphi$ 是 $n$ 维酉空间 $V$ 上的线性算子，证明：$\varphi$ 正规的充分必要条件是 $\varphi=\omega\psi$，其中 $\omega$ 为酉算子，$\psi$ 为半正定自伴随算子且 $\omega\psi=\psi\omega$.

3. 证明：谱分解定理之逆也成立，即若酉空间 $V$ 上存在一组线性算子 $\{E_1,E_2,\cdots,E_k\}$，适合 $E_1+E_2+\cdots+E_k=I$，$E_iE_j=O\ (i\ne j)$，$E_i^2=E_i=E_i^*$，并且线性算子 $\varphi=\lambda_1E_1+\lambda_2E_2+\cdots+\lambda_kE_k$，则 $\varphi$ 是正规算子.

4. 求证：$n$ 阶实对称阵 $A$ 是半正定阵的充分必要条件是存在同阶实对称阵 $B$，使 $A=B^2$.

5. 设 $A$ 为 $m\times n$ 列满秩实矩阵，$P=A(A'A)^{-1}A'$，求证：存在 $m\times n$ 实矩阵 $Q$，使 $I_n=Q'Q$ 且 $P=QQ'$.

6. 设 $A$ 为 $n$ 阶实矩阵，求证：存在可逆矩阵 $Q$，使 $QAQ=A'$.

7. 求下列非异阵的极分解：
\[
\begin{array}{lll}
(1)\ \left(\begin{array}{rr}11 & -5\\ -2 & 10\end{array}\right);
&
(2)\ \left(\begin{array}{rr}1 & 1\\ 2 & -2\end{array}\right);
&
(3)\ \left(\begin{array}{rrr}20 & 4 & 0\\ 0 & 0 & 1\\ 4 & 20 & 0\end{array}\right).
\end{array}
\]

\section*{* § 9.9\quad 奇异值分解}

\noindent\textbf{定义 9.9.1}\quad 设 $A$ 是 $m\times n$ 实矩阵，如果存在非负实数 $\sigma$ 以及 $n$ 维非零实列向量 $\alpha$，$m$ 维非零实列向量 $\beta$，使
\[
  A\alpha=\sigma\beta,\quad A'\beta=\sigma\alpha,
\]
则称 $\sigma$ 是 $A$ 的奇异值，$\alpha,\beta$ 分别称为 $A$ 关于 $\sigma$ 的右奇异向量与左奇异向量.

为了从几何上描述奇异值问题，我们引入线性映射的伴随概念，它可以看成是内积空间上线性变换的伴随概念的推广.

\noindent\textbf{定义 9.9.2}\quad 设 $V,U$ 分别是 $n$ 维，$m$ 维欧氏空间，$\varphi$ 是 $V\to U$ 的线性映射. 若存在 $U\to V$ 的线性映射 $\varphi^*$，使对任意的 $v\in V,u\in U$，都有
\[
  (\varphi(v),u)=(v,\varphi^*(u))
\]
成立，则称 $\varphi^*$ 是 $\varphi$ 的伴随.

\noindent\textbf{定理 9.9.1}\quad 设 $V,U$ 分别是 $n$ 维，$m$ 维欧氏空间，$\varphi$ 是 $V\to U$ 的线性映射，则 $\varphi$ 的伴随 $\varphi^*$ 存在且唯一.

\noindent\textbf{证明}\quad 与定理 9.3.1（即线性变换伴随的存在唯一性）的证明类似，故从略. $\square$

从伴随的定义我们不难发现，若取定 $V$ 的一组标准正交基 $\{e_1,e_2,\cdots,e_n\}$，$U$ 的一组标准正交基 $\{f_1,f_2,\cdots,f_m\}$，设 $\varphi$ 在这两组基下的表示矩阵为 $A$，则 $\varphi^*$ 在这两组基下的表示矩阵为 $A'$，证明也和线性变换的情形相同. 因此，奇异值与奇异向量的几何定义即为下列等式成立：
\[
  \varphi(v)=\sigma u,\quad \varphi^*(u)=\sigma v,
\]
其中 $\sigma\ge0$，$v\in V,u\in U$ 都是非零向量. 不难验证 $\varphi^*\varphi$ 是 $V$ 上的半正定自伴随算子，$\varphi\varphi^*$ 是 $U$ 上的半正定自伴随算子. 又
\[
  \varphi^*\varphi(v)=\varphi^*(\sigma u)=\sigma\varphi^*(u)=\sigma^2v,
\]
因此，$\sigma^2$ 是 $\varphi^*\varphi$ 的特征值，$v$ 是 $\varphi^*\varphi$ 的属于 $\sigma^2$ 的特征向量. 同理，$\sigma^2$ 也是 $\varphi\varphi^*$ 的特征值，$u$ 是 $\varphi\varphi^*$ 的属于 $\sigma^2$ 的特征向量.

\noindent\textbf{定理 9.9.2}\quad 设 $V,U$ 分别是 $n$ 维，$m$ 维欧氏空间，$\varphi$ 是 $V\to U$ 的线性映射，则存在 $V$ 和 $U$ 的标准正交基，使 $\varphi$ 在这两组基下的表示矩阵为
\[
  \left(\begin{array}{cc}
    S & O\\
    O & O
  \end{array}\right),
\]
其中
\[
  S=
  \left(\begin{array}{cccc}
    \sigma_1 & & &\\
    & \sigma_2 & &\\
    & & \ddots &\\
    & & & \sigma_r
  \end{array}\right)
\]
是一个 $r$ 阶对角阵，$\sigma_1\ge\sigma_2\ge\cdots\ge\sigma_r>0$ 是 $\varphi$ 的非零奇异值.

\noindent\textbf{证明}\quad 因为 $\varphi^*\varphi$ 是 $V$ 上的半正定自伴随算子，故存在 $V$ 的一组标准正交基 $\{e_1,e_2,\cdots,e_n\}$，使 $\varphi^*\varphi$ 在这组基下的表示矩阵为 $n$ 阶对角阵 $\operatorname{diag}\{\lambda_1,\cdots,\lambda_r,0,\cdots,0\}$，其中 $r=\operatorname{r}(\varphi^*\varphi)=\operatorname{r}(\varphi)$ 且 $\lambda_1\ge\cdots\ge\lambda_r>0$ 为 $\varphi^*\varphi$ 的正特征值，从而有
\[
  \varphi^*\varphi(e_i)=\lambda_ie_i,\quad 1\le i\le r;\quad
  \varphi^*\varphi(e_j)=0,\quad r+1\le j\le n.
\]
令 $\sigma_i=\sqrt{\lambda_i}\ (i=1,2,\cdots,r)$ 为算术平方根. 注意到对任意的 $1\le i\le r$，
\[
  \|\varphi(e_i)\|^2=(\varphi(e_i),\varphi(e_i))=(\varphi^*\varphi(e_i),e_i)=\lambda_i(e_i,e_i)=\sigma_i^2\|e_i\|^2=\sigma_i^2,
\]
即 $\|\varphi(e_i)\|=\sigma_i$；对任意的 $1\le i\ne j\le r$，
\[
  (\varphi(e_i),\varphi(e_j))=(\varphi^*\varphi(e_i),e_j)=\lambda_i(e_i,e_j)=0;
\]
又对任意的 $r+1\le j\le n$，
\[
  \|\varphi(e_j)\|^2=(\varphi(e_j),\varphi(e_j))=(\varphi^*\varphi(e_j),e_j)=0,
\]
即 $\varphi(e_j)=0$. 令
\[
  f_i=\frac1{\sigma_i}\varphi(e_i),\quad i=1,2,\cdots,r,
\]
则 $f_1,f_2,\cdots,f_r$ 是 $U$ 中一组两两正交长度为 $1$ 的 $m$ 维列向量，由定理 9.2.2 可将它们扩张为 $U$ 的一组标准正交基 $\{f_1,\cdots,f_r,f_{r+1},\cdots,f_m\}$. 于是在 $V$ 的标准正交基 $\{e_1,e_2,\cdots,e_n\}$ 和 $U$ 的标准正交基 $\{f_1,f_2,\cdots,f_m\}$ 下，$\varphi$ 满足：
\[
  \varphi(e_i)=\sigma_if_i,\quad 1\le i\le r;\quad \varphi(e_j)=0,\quad r+1\le j\le n.
\]
由 $\varphi^*$ 在上述两组标准正交基下的表示矩阵是 $\varphi$ 的表示矩阵的转置可得
\[
  \varphi^*(f_i)=\sigma_ie_i,\quad 1\le i\le r;\quad \varphi^*(f_j)=0,\quad r+1\le j\le m,
\]
这就得到了要证的结论. $\square$

\noindent\textbf{推论 9.9.1}\quad 设 $A$ 是 $m\times n$ 实矩阵，则存在 $m$ 阶正交矩阵 $P$ 以及 $n$ 阶正交矩阵 $Q$，使
\[
  P'AQ=
  \left(\begin{array}{cc}
    S & O\\
    O & O
  \end{array}\right),
\]
其中
\[
  S=
  \left(\begin{array}{cccc}
    \sigma_1 & & &\\
    & \sigma_2 & &\\
    & & \ddots &\\
    & & & \sigma_r
  \end{array}\right)
\]
是一个 $r$ 阶对角阵，$\sigma_1\ge\sigma_2\ge\cdots\ge\sigma_r>0$ 是 $A$ 的非零奇异值.

\[
  P'AQ=
  \left(\begin{array}{cc}
    S & O\\
    O & O
  \end{array}\right)
\]
称为矩阵 $A$ 的正交相抵标准型，而
\[
  A=P
  \left(\begin{array}{cc}
    S & O\\
    O & O
  \end{array}\right)
  Q'
\]
则称为矩阵 $A$ 的奇异值分解. 矩阵的奇异值分解在信息理论、控制理论和大数据科学等领域有着重要的应用.

如何计算 $m\times n$ 实矩阵 $A$ 的奇异值分解？事实上，从定理 9.9.2 的证明过程中可以得到具体的计算方法. 首先，求出 $A'A$ 的正交相似标准型，即求出 $n$ 阶正交矩阵 $Q$，使
\[
  Q'A'AQ=\operatorname{diag}\{\lambda_1,\cdots,\lambda_r,0,\cdots,0\},
\]
其中 $r=\operatorname{r}(A'A)=\operatorname{r}(A)$ 且 $\lambda_1\ge\cdots\ge\lambda_r>0$ 为 $A'A$ 的正特征值. 其次，设 $Q=(\alpha_1,\alpha_2,\cdots,\alpha_n)$ 为列分块，令
\[
  \sigma_i=\sqrt{\lambda_i},\quad \beta_i=\frac1{\sigma_i}A\alpha_i,\quad i=1,2,\cdots,r,
\]
则 $\beta_1,\beta_2,\cdots,\beta_r$ 是两两正交长度为 $1$ 的 $m$ 维列向量，将它们扩张为 $\mathbb R_m$ 的一组标准正交基 $\{\beta_1,\beta_2,\cdots,\beta_m\}$. 最后，令 $P=(\beta_1,\beta_2,\cdots,\beta_m)$ 为 $m$ 阶正交矩阵，则
\[
  AQ=P
  \left(\begin{array}{cc}
    S & O\\
    O & O
  \end{array}\right),
\]
从而
\begin{equation}
  A=P
  \left(\begin{array}{cc}
    S & O\\
    O & O
  \end{array}\right)Q'
  \tag{9.9.1}
\end{equation}
即为 $A$ 的奇异值分解.

在上述计算过程中，正交矩阵 $Q$ 的选取并不唯一；当 $Q$ 取定之后，若 $\operatorname{r}(A)=r<m$，则正交矩阵 $P$ 的选取也不唯一. 因此在奇异值分解中，除了 $\operatorname{diag}\{S,O\}$（即 $A$ 的奇异值）是由 $A$ 唯一确定之外，正交矩阵 $P,Q$ 的选取一般都不唯一.

从 $n$ 阶矩阵 $A$ 的奇异值分解很容易得到 $A$ 的极分解. 事实上，由 (9.9.1) 式可得
\begin{equation}
  A=(PQ')Q
  \left(\begin{array}{cc}
    S & O\\
    O & O
  \end{array}\right)Q',
  \tag{9.9.2}
\end{equation}
其中 $R=PQ'$ 是 $n$ 阶正交矩阵，$B=Q\operatorname{diag}\{S,O\}Q'$ 是 $n$ 阶半正定实对称阵，从而 $A=RB$ 即为 $A$ 的极分解. 通过奇异值分解来求极分解，在处理奇异阵时很有用.

\noindent\textbf{例 9.9.1}\quad 求下列矩阵的奇异值分解和极分解：
\[
  A=
  \left(\begin{array}{ccc}
    1 & 2 & 5\\
    1 & 0 & 1\\
    0 & 1 & 2
  \end{array}\right).
\]

\noindent\textbf{解}\quad 经计算可得
\[
  A'A=
  \left(\begin{array}{ccc}
    2 & 2 & 6\\
    2 & 5 & 12\\
    6 & 12 & 30
  \end{array}\right).
\]
采用与例 9.5.2 相同的计算方法可得正交矩阵
\[
  Q=
  \left(\begin{array}{ccc}
    \dfrac1{\sqrt{30}} & \dfrac2{\sqrt5} & \dfrac1{\sqrt6}\\
    \dfrac2{\sqrt{30}} & -\dfrac1{\sqrt5} & \dfrac2{\sqrt6}\\
    \dfrac5{\sqrt{30}} & 0 & -\dfrac1{\sqrt6}
  \end{array}\right),
\]
使 $Q'A'AQ=\operatorname{diag}\{36,1,0\}$. 设 $Q$ 的 $3$ 个列向量依次为 $\alpha_1,\alpha_2,\alpha_3$，令
\[
  \sigma_1=6,\quad \beta_1=\frac16A\alpha_1=\frac1{\sqrt{30}}(5,1,2)',
\]
\[
  \sigma_2=1,\quad \beta_2=A\alpha_2=\frac1{\sqrt5}(0,2,-1)'.
\]
添加单位向量 $\beta_3=\dfrac1{\sqrt6}(-1,1,2)'$，使 $\beta_1,\beta_2,\beta_3$ 成为 $\mathbb R_3$ 的一组标准正交基. 令
\[
  P=(\beta_1,\beta_2,\beta_3),
\]
则 $P$ 为正交矩阵. 由 (9.9.1) 式可得 $A$ 的奇异值分解为
\[
  A=
  \left(\begin{array}{ccc}
    \dfrac5{\sqrt{30}} & 0 & -\dfrac1{\sqrt6}\\
    \dfrac1{\sqrt{30}} & \dfrac2{\sqrt5} & \dfrac1{\sqrt6}\\
    \dfrac2{\sqrt{30}} & -\dfrac1{\sqrt5} & \dfrac2{\sqrt6}
  \end{array}\right)
  \left(\begin{array}{ccc}
    6 & &\\
    & 1 &\\
    & & 0
  \end{array}\right)
  \left(\begin{array}{ccc}
    \dfrac1{\sqrt{30}} & \dfrac2{\sqrt{30}} & \dfrac5{\sqrt{30}}\\
    \dfrac2{\sqrt5} & -\dfrac1{\sqrt5} & 0\\
    \dfrac1{\sqrt6} & \dfrac2{\sqrt6} & -\dfrac1{\sqrt6}
  \end{array}\right).
\]
由 (9.9.2) 式可得 $A$ 的极分解为
\[
  A=
  \left(\begin{array}{ccc}
    0 & 0 & 1\\
    1 & 0 & 0\\
    0 & 1 & 0
  \end{array}\right)
  \left(\begin{array}{ccc}
    1 & 0 & 1\\
    0 & 1 & 2\\
    1 & 2 & 5
  \end{array}\right).
\]
如果我们选取 $-\beta_3$ 与 $\beta_1,\beta_2$ 组成 $\mathbb R_3$ 的一组标准正交基，令 $P_1=(\beta_1,\beta_2,-\beta_3)$，则可得 $A$ 的另一种奇异值分解，由此诱导的 $A$ 的另一种极分解为
\[
  A=
  \left(\begin{array}{ccc}
    \dfrac13 & \dfrac23 & \dfrac23\\
    \dfrac23 & -\dfrac23 & \dfrac13\\
    -\dfrac23 & -\dfrac13 & \dfrac23
  \end{array}\right)
  \left(\begin{array}{ccc}
    1 & 0 & 1\\
    0 & 1 & 2\\
    1 & 2 & 5
  \end{array}\right).
\]

\noindent\textbf{例 9.9.2}\quad 设 $V,U$ 分别是 $n$ 维，$m$ 维欧氏空间，$\varphi,\psi$ 是 $V\to U$ 的线性映射，$\varphi^*$ 和 $\psi^*$ 分别是 $\varphi$ 和 $\psi$ 的伴随. 若 $\varphi^*\varphi=\psi^*\psi$，证明：存在 $U$ 上的正交变换 $\omega$，使 $\varphi=\omega\psi$.

\noindent\textbf{证明}\quad 我们将沿用定理 9.9.2 证明中的记号. 设在 $V$ 的标准正交基 $\{e_1,e_2,\cdots,e_n\}$ 下，$\varphi^*\varphi=\psi^*\psi$ 的表示矩阵为 $n$ 阶对角阵 $\operatorname{diag}\{\lambda_1,\cdots,\lambda_r,0,\cdots,0\}$，则由定理 9.9.2 的证明知存在 $U$ 的标准正交基 $\{f_1,f_2,\cdots,f_m\}$，使 $\varphi$ 在这两组基下的表示矩阵为分块对角阵 $\operatorname{diag}\{S,O\}$. 同理，存在 $U$ 的另一组标准正交基 $\{g_1,g_2,\cdots,g_m\}$，使 $\psi$ 在 $\{e_1,e_2,\cdots,e_n\}$ 和 $\{g_1,g_2,\cdots,g_m\}$ 下的表示矩阵也是 $\operatorname{diag}\{S,O\}$. 现定义 $U$ 上的线性变换 $\omega$ 为 $\omega(g_i)=f_i\ (i=1,2,\cdots,m)$，则 $\omega$ 是 $U$ 上的正交变换且
\[
  \omega\psi(e_i)=\omega(\sigma_ig_i)=\sigma_i\omega(g_i)=\sigma_if_i=\varphi(e_i),\quad i=1,\cdots,r;
\]
\[
  \omega\psi(e_j)=\omega(0)=0=\varphi(e_j),\quad j=r+1,\cdots,n,
\]
故 $\varphi=\omega\psi$. $\square$

实矩阵（欧氏空间之间的线性映射）的奇异值分解理论完全可以类似地推广到复矩阵（酉空间之间的线性映射）的情形，我们把这些推广留给读者自己完成.

\section*{习题 9.9}

1. 求下列矩阵的奇异值分解：
\[
\begin{array}{lll}
(1)\ \left(\begin{array}{rr}1 & 1\\ 1 & 1\\ -1 & -1\end{array}\right);
&
(2)\ \left(\begin{array}{rr}1 & 1\\ 0 & 1\\ 1 & 0\\ 1 & 1\end{array}\right);
&
(3)\ \left(\begin{array}{rrr}1 & 1 & 1\\ 1 & 0 & -1\\ 1 & -2 & 1\\ 1 & 1 & 1\end{array}\right).
\end{array}
\]

2. 求下列矩阵的极分解：
\[
\begin{array}{ll}
(1)\ \left(\begin{array}{cc}1 & 2\\ 1 & 2\end{array}\right);
&
(2)\ \left(\begin{array}{rrrr}
1 & 1 & -3 & 1\\
-3 & 1 & 1 & 1\\
1 & 1 & 1 & -3\\
1 & -3 & 1 & 1
\end{array}\right).
\end{array}
\]

\section*{* § 9.10\quad 最小二乘解}

我们先来讨论一个几何问题. 设在三维欧氏空间 $V$ 中有一个平面 $W$，平面外有一点 $C$，现在要求这一点到该平面的最短距离. 为方便起见，不妨设该平面由 $V$ 的一组标准正交基 $\{e_1,e_2,e_3\}$ 中的两个向量 $e_1,e_2$ 张成（图 9.1）.

$C$ 为向量 $v$ 的终点，则 $v$ 到平面 $W$ 的距离 $d$ 等于 $C$ 到垂足点 $P$ 的距离，即
\[
  d=\|v-u\|.
\]
因此只需求出 $u$，就可以求出 $d$ 来. 由于 $u\in W$，故可设
\[
  u=\lambda_1e_1+\lambda_2e_2,
\]
由 $(v-u)\perp e_i\ (i=1,2)$，得
\[
  (v-u,e_i)=0,\quad i=1,2.
\]
因此
\[
  \lambda_1=(u,e_1)=(v,e_1),\quad \lambda_2=(u,e_2)=(v,e_2),
\]
即
\[
  u=(v,e_1)e_1+(v,e_2)e_2.
\]
另一方面，
\[
  v=(v,e_1)e_1+(v,e_2)e_2+(v,e_3)e_3,
\]
即有
\begin{equation}
  d=|(v,e_3)|.
  \tag{9.10.1}
\end{equation}
我们称向量 $u$ 为向量 $v$ 在子空间 $W$ 上的正交投影向量.

现在我们把问题提得更一般些. 设 $V$ 是一个内积空间，$W$ 是 $V$ 的子空间（$W\ne V$），$v$ 是 $V$ 中的一个向量，现要在 $W$ 中找一个向量 $u$，使 $\|v-u\|$ 最小. 如果这样的 $u$ 存在，则称 $u$ 是 $v$ 在 $W$ 中的正交投影向量，长度 $\|v-u\|$ 被称为 $v$ 到 $W$ 的距离.

\noindent\textbf{定理 9.10.1}\quad 设 $W$ 是有限维内积空间 $V$ 的子空间，$v\in V$，则

(1) 在 $W$ 中存在唯一的向量 $u$，使 $\|v-u\|$ 最小且这时 $(v-u)\perp W$；

(2) 若 $\{e_1,\cdots,e_m\}$ 是 $W$ 的标准正交基，又 $\{e_{m+1},\cdots,e_n\}$ 是 $W^\perp$ 的标准正交基，这样 $\{e_1,e_2,\cdots,e_n\}$ 就成为 $V$ 的一组标准正交基，则
\[
  u=(v,e_1)e_1+(v,e_2)e_2+\cdots+(v,e_m)e_m,
\]
\[
  v-u=(v,e_{m+1})e_{m+1}+\cdots+(v,e_n)e_n,
\]
\begin{equation}
  \|v-u\|=\left(|(v,e_{m+1})|^2+\cdots+|(v,e_n)|^2\right)^{\frac12}.
  \tag{9.10.2}
\end{equation}

\noindent\textbf{证明}\quad 设 $w$ 是 $W$ 中任一向量，要证明 $\|v-u\|\le\|v-w\|$. 注意到这时 $v-w=(v-u)+(u-w)$，其中 $(v-u)\perp W,u-w\in W$. 因此，由勾股定理有
\[
  \|v-w\|^2=\|v-u\|^2+\|u-w\|^2.
\]
若 $w\ne u$，则
\[
  \|v-w\|^2>\|v-u\|^2.
\]
这样我们不仅证明了 $\|v-u\|$ 的最小性，也证明了 $u$ 的唯一性. $\square$

现在我们举例来说明定理 9.10.1 在求最小二乘解方面的应用. 在许多实际问题中我们经常会碰到所谓的“矛盾线性方程组”的求解问题. 为说得更清楚一些，举一个简单的例子.

在经济学中，个人的收入与消费之间存在着密切的关系. 收入越多，消费水平也越高，收入较少，消费水平也较低. 从一个社会整体来看，个人的平均收入与平均消费之间大致呈线性关系. 若 $u$ 表示收入，$v$ 表示支出，则 $u,v$ 适合
\begin{equation}
  u=a+bv,
  \tag{9.10.3}
\end{equation}
其中 $a,b$ 是两个常数，需要根据具体的统计数据来确定. 假设现在有一组统计数据表示 $3$ 年中每年的收入与消费情况（表 9.1），现要根据这一组统计数据求出 $a,b$.

将 $u,v$ 的值代入 (9.10.3) 式得到一个两个未知数 $3$ 个方程式的线性方程组：
\[
  \begin{cases}
    a+1.2b=1.6,\\
    a+1.4b=1.7,\\
    a+1.8b=2.0.
  \end{cases}
\]
从第一、第二个方程式可求出 $a=1,b=0.5$，代入第三个方程式：
\[
  1+1.8\times0.5=1.9\ne2.0,
\]
这说明上面的线性方程组无解，那么这样一来是不是说我们的问题就没有意义了呢？当然不是！事实上，收入与消费的关系通常极为复杂，我们把它当成线性关系只是一种近似的假设. 另外，统计数据本身不可避免地会产生误差，也就是说统计表只是实际情况的近似反映. 我们的目的是求出 $a,b$ 的值以供理论分析之用. 既然统计数据有误差，就不可能也没有必要求出 $a,b$ 的精确解. 此矛盾的出现并不意味着我们因此而束手无策了. 我们可以对 $a,b$ 提出这样的要求：求出 $a$ 与 $b$，使得到的关系式 $u=a+bv$ 能尽可能好地符合实际情形. 用数学语言来说就是求 $a,b$，使平方偏差
\[
  \big((a+1.2b)-1.6\big)^2+\big((a+1.4b)-1.7\big)^2+\big((a+1.8b)-2.0\big)^2
\]
取最小值. 这就是所谓的最小二乘问题. 一般来说，为了使理论关系式更符合实际，通常要求统计数据多一点，即方程式的个数多一点.

如何求“最小二乘解”？让我们来作一些理论上的分析. 假设有下列矛盾线性方程组：
\[
  \begin{cases}
    a_{11}x_1+a_{12}x_2+\cdots+a_{1n}x_n=b_1,\\
    a_{21}x_1+a_{22}x_2+\cdots+a_{2n}x_n=b_2,\\
    \hspace{5em}\cdots\cdots\cdots\cdots\\
    a_{m1}x_1+a_{m2}x_2+\cdots+a_{mn}x_n=b_m,
  \end{cases}
\]
其中 $m>n$，它的系数矩阵记为 $A$. 为方便，将上述线性方程组写为矩阵形式：
\begin{equation}
  Ax=\beta,
  \tag{9.10.4}
\end{equation}
其中
\[
  x=
  \left(\begin{array}{c}
    x_1\\ x_2\\ \vdots\\ x_n
  \end{array}\right),\quad
  \beta=
  \left(\begin{array}{c}
    b_1\\ b_2\\ \vdots\\ b_m
  \end{array}\right).
\]
现要求出 $x$，使
\[
  \sum_{i=1}^m\left(\sum_{j=1}^na_{ij}x_j-b_i\right)^2
\]
取最小值. 记
\[
  \alpha_1=
  \left(\begin{array}{c}
    a_{11}\\ a_{21}\\ \vdots\\ a_{m1}
  \end{array}\right),\quad
  \alpha_2=
  \left(\begin{array}{c}
    a_{12}\\ a_{22}\\ \vdots\\ a_{m2}
  \end{array}\right),\quad
  \cdots,\quad
  \alpha_n=
  \left(\begin{array}{c}
    a_{1n}\\ a_{2n}\\ \vdots\\ a_{mn}
  \end{array}\right),
\]
考虑向量
\[
  x_1\alpha_1+x_2\alpha_2+\cdots+x_n\alpha_n-\beta,
\]
这是一个 $m$ 维列向量，第 $i$ 个坐标为 $\sum_{j=1}^na_{ij}x_j-b_i$. 因此，若记 $V=\mathbb R_m$，内积取标准内积，则
\[
  \|(x_1\alpha_1+x_2\alpha_2+\cdots+x_n\alpha_n)-\beta\|^2
  =\sum_{i=1}^m\left(\sum_{j=1}^na_{ij}x_j-b_i\right)^2.
\]
而
\[
  \|(x_1\alpha_1+x_2\alpha_2+\cdots+x_n\alpha_n)-\beta\|
  =\|Ax-\beta\|=\|\beta-Ax\|.
\]
当 $x_1,x_2,\cdots,x_n$ 变动时，$x_1\alpha_1+x_2\alpha_2+\cdots+x_n\alpha_n$ 就是由 $\alpha_1,\alpha_2,\cdots,\alpha_n$ 张成的 $V$ 的子空间，记为 $W$. 要使 $\|\beta-Ax\|$ 取最小值，实际上就是要求 $\beta$ 到 $W$ 的距离. 由定理 9.10.1 知道，只需取 $\beta$ 在 $W$ 上的正交投影向量 $\gamma$ 就可以了. 于是求方程组 (9.10.4) 的最小二乘解又归结为求下列线性方程组的解：
\begin{equation}
  Ax=\gamma.
  \tag{9.10.5}
\end{equation}
由定理 9.10.1 知道这样的解总是存在的. 我们可以按照定理 9.10.1 中的办法先求 $W$ 的标准正交基，再扩展为 $V$ 的标准正交基，从而求出 $\gamma$，最后解线性方程组 (9.10.5) 便可求出 $x$. 但是这样做比较麻烦，我们希望能有一个比较简便的方法求出 $x$ 来.

首先我们注意到在实际问题中，矩阵 $A$ 的秩通常都等于 $n$. 因此，可设 $A$ 是列满秩阵，即 $\operatorname{r}(A)=n$. 这就是说，向量 $\alpha_1,\alpha_2,\cdots,\alpha_n$ 线性无关，$W$ 是由 $\alpha_1,\alpha_2,\cdots,\alpha_n$ 张成的 $n$ 维子空间. 设
\begin{equation}
  \gamma=x_1\alpha_1+x_2\alpha_2+\cdots+x_n\alpha_n,
  \tag{9.10.6}
\end{equation}
由定理 9.10.1 知道 $(\beta-\gamma)\perp W$，因此 $(\beta-\gamma)\perp\alpha_i\ (i=1,2,\cdots,n)$，即 $((\beta-\gamma),\alpha_i)=0$，或者 $(\gamma,\alpha_i)=(\beta,\alpha_i)$. 这样便有下列线性方程组：
\[
  \begin{cases}
    (\alpha_1,\alpha_1)x_1+(\alpha_2,\alpha_1)x_2+\cdots+(\alpha_n,\alpha_1)x_n=(\beta,\alpha_1),\\
    (\alpha_1,\alpha_2)x_1+(\alpha_2,\alpha_2)x_2+\cdots+(\alpha_n,\alpha_2)x_n=(\beta,\alpha_2),\\
    \hspace{5em}\cdots\cdots\cdots\cdots\\
    (\alpha_1,\alpha_n)x_1+(\alpha_2,\alpha_n)x_2+\cdots+(\alpha_n,\alpha_n)x_n=(\beta,\alpha_n).
  \end{cases}
\]
若记 $A'$ 为 $A$ 的转置，则上面的方程组可写为矩阵形式：
\[
  A'Ax=A'\beta.
\]
由于 $A$ 的秩为 $n$，故 $A'A$ 的秩也为 $n$，即 $A'A$ 为 $n$ 阶非异阵，于是
\[
  x=(A'A)^{-1}A'\beta.
\]
这就是线性方程组 (9.10.4) 的最小二乘解.

下面我们来求前面例子中提出的线性方程组的最小二乘解：
\[
  \begin{cases}
    a+1.2b=1.6,\\
    a+1.4b=1.7,\\
    a+1.8b=2.0.
  \end{cases}
\]
这时
\[
  A=
  \left(\begin{array}{cc}
    1 & 1.2\\
    1 & 1.4\\
    1 & 1.8
  \end{array}\right),\quad
  A'=
  \left(\begin{array}{ccc}
    1 & 1 & 1\\
    1.2 & 1.4 & 1.8
  \end{array}\right),\quad
  \beta=
  \left(\begin{array}{c}
    1.6\\ 1.7\\ 2.0
  \end{array}\right).
\]
不难看出 $A$ 的秩等于 $2$，并通过计算得
\[
  A'A=
  \left(\begin{array}{cc}
    3 & 4.4\\
    4.4 & 6.64
  \end{array}\right),\quad
  (A'A)^{-1}=\frac1{0.56}
  \left(\begin{array}{rr}
    6.64 & -4.4\\
    -4.4 & 3
  \end{array}\right),
\]
\[
  A'\beta=
  \left(\begin{array}{c}
    5.3\\
    7.9
  \end{array}\right),\quad
  (A'A)^{-1}A'\beta=\frac1{0.56}
  \left(\begin{array}{c}
    0.432\\
    0.38
  \end{array}\right)
  \approx
  \left(\begin{array}{c}
    0.77\\
    0.68
  \end{array}\right).
\]
由此得最小二乘解（近似值）：
\[
  \begin{cases}
    a=0.77,\\
    b=0.68.
  \end{cases}
\]
因此，收入与消费的关系式为
\[
  u=0.77+0.68v.
\]

\section*{习题 9.10}

1. 设 $A$ 为 $m\times n$ 实矩阵，求证：对任意的 $m$ 维实列向量 $\beta$，$n$ 元线性方程组 $A'Ax=A'\beta$ 一定有解.

\end{document}
